\chapter{Recursion}
\chapterepigraph{To understand recursion, trust the recursion. Assume
your function already correctly solves every strictly smaller version of
the problem, and ask only one question: given those correct smaller
answers, how do I build the answer to \emph{this} input? If you can
answer that question once, you have solved every input, by induction.}

\section{The One Question Every Recursive Function Must Answer}
\label{sec:rec-intro}

Every recursive function is built from exactly three ingredients: a
\textbf{base case} (the smallest input(s) the function can answer
directly, without recursing), a \textbf{recursive case} (how to build
this input's answer from one or more strictly smaller sub-calls, trusted
to already be correct), and a guarantee that the sub-calls' inputs
actually \emph{shrink} toward the base case (otherwise the recursion
never terminates).

\begin{ideabox}[The Recursive Leap of Faith]
When designing a recursive function, do not try to mentally trace through
every recursive call -- this is the single most common source of
confusion and wasted time. Instead, write down what the function is
\emph{supposed} to do in words, assume every call with a strictly smaller
input already does exactly that correctly, and then write only the logic
needed to combine those trusted smaller answers into this input's answer.
Verifying the base case and the combination step is enough to prove the
whole function correct, by mathematical induction on the input's size --
there is no need to simulate the call stack by hand except when actively
debugging.
\end{ideabox}

\subsection*{Two Families of Recursion in This Chapter}

\begin{patternbox}[How to Recognise Which Family You Are In]
\begin{itemize}
  \item \textbf{``Compute a single value/result''} (a power, a digit
        count, a stack reversed or sorted) -- \textbf{straight-line
        recursion}: each call makes exactly one (or a small fixed
        number of) recursive call(s) and combines the result(s) in a
        fixed way. The skill here is finding the right way to shrink the
        input and the right combination step, often exploiting a
        mathematical identity (as in fast exponentiation) to shrink
        faster than one unit at a time.
  \item \textbf{``Generate/find all ...''} (subsequences, subsets,
        combinations, permutations, paths, valid placements) --
        \textbf{exploration recursion}, commonly called
        \textbf{backtracking}: each call branches into \emph{multiple}
        recursive calls, one per choice available at this step, and an
        explicit \emph{undo} step restores shared state after each
        branch returns, so that the next branch starts from a clean
        slate. The skill here is identifying, at each step, the full set
        of choices available, and recognizing when a partial choice can
        be abandoned early (pruned) because it can no longer lead to a
        valid complete answer.
\end{itemize}
This chapter moves from the first family (Section~\ref{sec:rec-good-numbers}
through Section~\ref{sec:rec-sort-stack}) into the second
(Section~\ref{sec:rec-power-set} onward), building backtracking's
``choose, explore, un-choose'' template up gradually through string and
subset generation before tackling full constraint-satisfaction problems
like N-Queens and Sudoku.
\end{patternbox}

\subsection*{Why Recursion Sets Up Everything in the Dynamic Programming Chapters}

Every dynamic programming problem in the chapters that follow
\emph{begins} as a recursive solution: identify the choices available at
each step, write the recursive recurrence that combines smaller
sub-answers into this one, and verify it against small examples by hand
-- exactly the discipline built in this chapter. Memoization and
tabulation, covered there, are mechanical transformations applied
\emph{after} a correct recursive solution already exists; they do not
replace the need to find that recursion first. Mastering straight-line
and exploration recursion here is therefore not a detour before dynamic
programming -- it is dynamic programming's first and most important step,
practiced in isolation.

Each section in this chapter follows the established structure: problem
statement and keywords; the reasoning that identifies the base case,
the recursive case, and (for backtracking problems) the choice set and
any pruning opportunity; the algorithm in words and pseudocode; a
hand-traced dry run (often shown as a small recursion tree); complete
compilable C++ code; and a time/space complexity analysis, including the
recursion's call-stack depth.

\newpage

\section{Count Good Numbers}
\label{sec:rec-good-numbers}

\problemstatement{A digit string of length $n$ is called \textbf{good} if
every digit at an \textbf{even} index (0-indexed) is one of the even
digits $\{0,2,4,6,8\}$, and every digit at an \textbf{odd} index is one
of the prime digits $\{2,3,5,7\}$. Count how many good digit strings of
length $n$ exist, modulo $10^9+7$.}

\begin{patternbox}
Despite involving ``counting all possibilities,'' this is \emph{not} a
backtracking problem -- there is no need to actually generate any digit
string, because each position's choice is completely \textbf{independent}
of every other position's choice. Whenever every position in a sequence
can be filled independently from a fixed-size choice set, the total count
is simply the \textbf{product} of each position's choice count, and that
product is best computed via \textbf{fast exponentiation}, a
straight-line recursion that shrinks an exponent by half at each step.
\end{patternbox}

\subsection*{Understanding the problem carefully}

Among the $n$ positions, $\lceil n/2 \rceil$ are at even indices ($0, 2,
4, \dots$) and $\lfloor n/2 \rfloor$ are at odd indices. Each even-index
position independently has $5$ valid choices ($0,2,4,6,8$), and each
odd-index position independently has $4$ valid choices ($2,3,5,7$). By
the multiplication principle (every combination of independent choices
is a distinct valid string), the total count is
\[
  5^{\lceil n/2 \rceil} \times 4^{\lfloor n/2 \rfloor} \pmod{10^9+7}.
\]
The only remaining challenge is computing a large power modulo $10^9+7$
efficiently, since $n$ can be large enough that computing $5^{n/2}$ by
repeated multiplication, one factor at a time, would be too slow.

\begin{ideabox}[Fast Exponentiation: Halve the Exponent, Not Decrement It]
To compute $\textit{base}^{\textit{exp}}$: if $\textit{exp}=0$, the
answer is $1$ (base case). If $\textit{exp}$ is even, recursively compute
$\textit{half} = \textit{base}^{\textit{exp}/2}$, and the answer is
$\textit{half} \times \textit{half}$. If $\textit{exp}$ is odd, compute
$\textit{half} = \textit{base}^{(\textit{exp}-1)/2}$, and the answer is
$\textit{half} \times \textit{half} \times \textit{base}$.
\end{ideabox}

\subsection*{Why halving the exponent gives $O(\log n)$ instead of $O(n)$}

A naive recursive power function that decrements the exponent by $1$ at
each call would need $O(n)$ recursive calls, each doing one
multiplication. Halving the exponent instead means the recursion depth is
only $O(\log n)$ -- each level squares an already-computed result rather
than multiplying by the base value just once, exploiting the identity
$\textit{base}^{2k} = (\textit{base}^k)^2$ to skip directly from
computing one power to a power roughly twice as large, in a single
multiplication. This is the same ``shrink faster than one unit at a
time'' idea elaborated fully in Section~\ref{sec:rec-pow}.

\subsection*{Algorithm}

\begin{enumerate}
  \item $\textit{evenCount} \gets \lceil n/2 \rceil$,
        $\textit{oddCount} \gets \lfloor n/2 \rfloor$.
  \item Compute $\textit{a} \gets 5^{\textit{evenCount}} \bmod
        (10^9+7)$ via fast exponentiation.
  \item Compute $\textit{b} \gets 4^{\textit{oddCount}} \bmod
        (10^9+7)$ via fast exponentiation.
  \item Return $(\textit{a} \times \textit{b}) \bmod (10^9+7)$.
\end{enumerate}

\subsection*{Dry run}

\dryrun{Take $n=4$.}

$\textit{evenCount}=2$ (indices $0,2$), $\textit{oddCount}=2$ (indices
$1,3$). $5^2=25$, $4^2=16$. Answer $=25\times16=400$.

Checking by hand for a tiny case, $n=1$: $\textit{evenCount}=1$,
$\textit{oddCount}=0$. $5^1=5$, $4^0=1$. Answer $=5$ -- correctly, the
$5$ single-digit good numbers are $0,2,4,6,8$ (index $0$ is even, so any
even digit qualifies, and there is no odd index to satisfy).

\subsection*{C++ implementation}

\begin{lstlisting}
#include <bits/stdc++.h>
using namespace std;
const long long MOD = 1e9 + 7;

long long power(long long base, long long exp) {
    if (exp == 0) return 1;
    long long half = power(base, exp / 2) % MOD;
    long long result = (half * half) % MOD;
    if (exp % 2 == 1) result = (result * base) % MOD;
    return result;
}

int countGoodNumbers(long long n) {
    long long evenCount = (n + 1) / 2;
    long long oddCount = n / 2;

    long long a = power(5, evenCount);
    long long b = power(4, oddCount);
    return (a * b) % MOD;
}

int main() {
    cout << "Good numbers of length 4: " << countGoodNumbers(4) << endl;
    return 0;
}
\end{lstlisting}

\begin{complexitybox}
\textbf{Time Complexity.} $O(\log n)$ for each fast-exponentiation call,
giving an overall time complexity of
\[
  O(\log n).
\]
\textbf{Space Complexity.} $O(\log n)$ for the recursion stack of fast
exponentiation.
\end{complexitybox}

\begin{takeawaybox}
Whenever a counting problem's choices at every position are independent
of each other, resist the urge to enumerate -- multiply the per-position
choice counts directly, and reach for fast exponentiation the moment that
product involves a single base raised to a large power. This pattern of
recognizing independence to avoid generation entirely recurs throughout
combinatorial counting problems.
\end{takeawaybox}

\section{Implement Pow(x, n)}
\label{sec:rec-pow}

\problemstatement{Implement a function that computes $x^n$ for a real
number $x$ and an integer $n$, which may be negative, in better than
$O(|n|)$ time.}

\begin{patternbox}
\emph{``Compute $x^n$ efficiently''} is the canonical introduction to
\textbf{fast exponentiation} (also called binary exponentiation), already
used as a building block in Section~\ref{sec:rec-good-numbers}. The
keyword to internalize generally: whenever a recursive shrinking step can
\textbf{halve} the input rather than merely decrement it by one, the
recursion depth drops from $O(n)$ to $O(\log n)$ -- a pattern that
reappears constantly, not just for exponentiation.
\end{patternbox}

\subsection*{Understanding the problem carefully}

The naive recursive definition $x^n = x \times x^{n-1}$ (with base case
$x^0=1$) is correct but requires $n$ recursive calls. The key
algebraic identity that lets us do better is
\[
  x^n = \left(x^{n/2}\right)^2 \quad \text{when } n \text{ is even}, \qquad
  x^n = x \times \left(x^{(n-1)/2}\right)^2 \quad \text{when } n \text{ is
  odd}.
\]
Either way, computing $x^n$ only requires computing \emph{one} smaller
power, $x^{\lfloor n/2 \rfloor}$, and then a constant amount of extra
multiplication -- the exponent is halved at every recursive step instead
of merely decremented.

\begin{ideabox}[Halve the Exponent, Handle Negative Exponents Separately]
If $n=0$: return $1$ (base case). If $n<0$: compute $x^{-n}$ using the
same recursion, and return its reciprocal $1/x^{-n}$ at the very end (do
not try to handle negative exponents inside the recursive halving step
itself -- convert once, up front, and recurse only on non-negative
exponents). For $n>0$: recursively compute $\textit{half} =
x^{\lfloor n/2 \rfloor}$; if $n$ is even, return $\textit{half} \times
\textit{half}$; if $n$ is odd, return $\textit{half} \times \textit{half}
\times x$.
\end{ideabox}

\subsection*{Why halving gives logarithmic depth}

Each recursive call's exponent is at most half of its caller's exponent,
so after $k$ calls the exponent has shrunk by a factor of $2^k$ -- it
reaches the base case $0$ once $2^k \ge n$, i.e.\ after only
$O(\log n)$ calls, regardless of how large $n$ is. This is in sharp
contrast to the naive decrement-by-one recursion, whose depth is exactly
$n$. The trade-off is that each call here does a small constant amount
of extra multiplication work (squaring the smaller result, and possibly
one more multiplication by $x$ for odd exponents), but this constant
overhead is negligible compared to the exponential reduction in the
number of calls needed.

\subsection*{Algorithm}

\begin{enumerate}
  \item If $n<0$: return $1 / \textsc{Pow}(x, -n)$.
  \item $\textsc{Pow}(x, n)$:
  \begin{enumerate}
    \item If $n=0$: return $1$.
    \item $\textit{half} \gets \textsc{Pow}(x, n/2)$ (integer division).
    \item If $n$ is even: return $\textit{half} \times \textit{half}$.
    \item Else: return $\textit{half} \times \textit{half} \times x$.
  \end{enumerate}
\end{enumerate}

\subsection*{Dry run}

\dryrun{Compute $2^{10}$.}

\begin{itemize}
  \item $\textsc{Pow}(2,10)$: $10$ even. $\textit{half}=\textsc{Pow}(2,5)$.
  \item $\textsc{Pow}(2,5)$: $5$ odd. $\textit{half}=\textsc{Pow}(2,2)$.
  \item $\textsc{Pow}(2,2)$: $2$ even. $\textit{half}=\textsc{Pow}(2,1)$.
  \item $\textsc{Pow}(2,1)$: $1$ odd. $\textit{half}=\textsc{Pow}(2,0)=1$.
        Return $1\times1\times2=2$.
  \item Back at $\textsc{Pow}(2,2)$: $\textit{half}=2$. Return
        $2\times2=4$.
  \item Back at $\textsc{Pow}(2,5)$: $\textit{half}=4$. Return
        $4\times4\times2=32$.
  \item Back at $\textsc{Pow}(2,10)$: $\textit{half}=32$. Return
        $32\times32=1024$.
\end{itemize}

$2^{10}=1024$, confirmed. Only $4$ recursive calls were needed for an
exponent of $10$, instead of $10$ calls a naive decrement-by-one
recursion would have required.

\subsection*{C++ implementation}

\begin{lstlisting}
#include <bits/stdc++.h>
using namespace std;

double fastPow(double x, long long n) {
    if (n == 0) return 1.0;
    double half = fastPow(x, n / 2);
    double result = half * half;
    if (n % 2 == 1) result *= x;
    return result;
}

double myPow(double x, int n) {
    long long N = n;
    if (N < 0) return 1.0 / fastPow(x, -N);
    return fastPow(x, N);
}

int main() {
    cout << myPow(2.0, 10) << endl;   // 1024
    cout << myPow(2.0, -2) << endl;   // 0.25
    return 0;
}
\end{lstlisting}

\begin{complexitybox}
\textbf{Time Complexity.}
\[
  O(\log n).
\]
\textbf{Space Complexity.} $O(\log n)$ for the recursion stack (or
$O(1)$ if rewritten as an iterative loop that repeatedly squares the
base and halves the exponent, accumulating a result whenever the current
bit of the exponent is set).
\end{complexitybox}

\begin{takeawaybox}
Fast exponentiation is the prototype for an entire family of techniques:
whenever an operation is associative (like multiplication) and you need
to apply it $n$ times, check whether the operation can be ``doubled up''
(squared, in this case) to cut the required number of steps from $O(n)$
to $O(\log n)$. The same halving idea underlies fast matrix
exponentiation for linear recurrences and binary lifting for tree
ancestor queries, both common extensions of this exact pattern.
\end{takeawaybox}

\section{Recursive Implementation of atoi}
\label{sec:rec-atoi}

\problemstatement{Implement (recursively) the C-style $\textit{atoi}$
function: given a string, skip any leading whitespace, then an optional
$\texttt{+}$ or $\texttt{-}$ sign, then consume as many leading digits as
possible, and return the resulting integer -- clamped to the 32-bit
signed integer range $[-2^{31}, 2^{31}-1]$ if it would otherwise
overflow. If no valid number can be parsed at all, return $0$.}

\begin{patternbox}
This problem has no clever algorithmic trick -- the difficulty is purely
in correctly handling several small edge cases (whitespace, sign,
leading zeros, early non-digit characters, overflow) using recursion
instead of an iterative loop. The keyword to notice: \emph{``recursive
implementation of''} an otherwise simple iterative task is an exercise in
expressing a left-to-right scan as a chain of recursive calls, each
processing one character and delegating the rest to a sub-call.
\end{patternbox}

\subsection*{Understanding the problem carefully}

The parsing naturally splits into three independent phases, each best
handled by its own small recursive (or simple iterative) helper: skipping
whitespace, consuming an optional sign, and then accumulating digits one
at a time. The digit-accumulation phase is the one most naturally
expressed recursively: at each step, if the current character is a
digit, fold it into a running accumulated value and recurse on the rest
of the string; the moment a non-digit (or the string's end) is reached,
stop and return the accumulated value so far.

\begin{ideabox}[An Accumulator Passed Down Through the Recursion]
$\textsc{ParseDigits}(s, i, \textit{acc})$: if $i$ is out of bounds or
$s[i]$ is not a digit, return $\textit{acc}$ (base case -- nothing more
to consume). Otherwise, compute $\textit{newAcc} = \textit{acc} \times 10
+ (s[i] - \texttt{'0'})$, clamp $\textit{newAcc}$ immediately if it has
already exceeded the representable range (to avoid undefined-behavior
overflow in the accumulator itself), and recurse:
$\textsc{ParseDigits}(s, i+1, \textit{newAcc})$.
\end{ideabox}

\subsection*{Why clamping during accumulation, not just at the end, matters}

If the accumulator is allowed to grow without bound while scanning a
string of many digits (e.g.\ a string of $50$ digits), it can overflow
even a 64-bit accumulator long before the recursion finishes, leading to
incorrect (wrapped-around) results rather than a clean clamp. Checking
and clamping the accumulator the moment it \emph{exceeds} the valid
32-bit range -- and, once clamped, simply continuing to consume (and
ignore) any remaining digits so the function still terminates at the
string's actual end of digits -- avoids this, ensuring the final
returned value is always either the true parsed integer or correctly
clamped to $\pm 2^{31}{-}1$ /$-2^{31}$.

\subsection*{Algorithm}

\begin{enumerate}
  \item Skip leading whitespace characters (advance an index $i$ while
        $s[i]$ is a space).
  \item If $s[i]$ is $\texttt{+}$ or $\texttt{-}$: record the sign,
        advance $i$.
  \item Call $\textsc{ParseDigits}(s, i, 0)$ to accumulate the digit
        value, clamping along the way as described above.
  \item Apply the recorded sign to the result, then clamp the
        signed result to $[-2^{31}, 2^{31}-1]$ one final time.
\end{enumerate}

\subsection*{Dry run}

\dryrun{Take $s = \texttt{"  -042abc"}$.}

\begin{itemize}
  \item Skip whitespace: $i$ advances past the two leading spaces.
  \item Sign: $s[i]=\texttt{'-'}$, record negative, advance $i$.
  \item $\textsc{ParseDigits}$: consume \texttt{0} ($\textit{acc}=0$),
        \texttt{4} ($\textit{acc}=4$), \texttt{2} ($\textit{acc}=42$);
        next character is \texttt{a}, not a digit -- stop, return $42$.
  \item Apply sign: $-42$. Well within range, no clamping needed.
\end{itemize}

Result: $-42$.

\subsection*{C++ implementation}

\begin{lstlisting}
#include <bits/stdc++.h>
using namespace std;

long long parseDigits(const string& s, int i, long long acc) {
    if (i >= (int)s.size() || !isdigit(s[i])) return acc;

    long long newAcc = acc * 10 + (s[i] - '0');
    if (newAcc > INT_MAX) newAcc = (long long)INT_MAX + 1; // clamp marker
    return parseDigits(s, i + 1, newAcc);
}

int myAtoi(string s) {
    int i = 0, n = s.size();
    while (i < n && s[i] == ' ') i++;

    int sign = 1;
    if (i < n && (s[i] == '+' || s[i] == '-')) {
        if (s[i] == '-') sign = -1;
        i++;
    }

    long long result = parseDigits(s, i, 0) * sign;

    if (result > INT_MAX) return INT_MAX;
    if (result < INT_MIN) return INT_MIN;
    return (int)result;
}

int main() {
    cout << myAtoi("  -042abc") << endl;          // -42
    cout << myAtoi("91283472332") << endl;        // 2147483647 (clamped)
    return 0;
}
\end{lstlisting}

\begin{complexitybox}
\textbf{Time Complexity.} $O(L)$, where $L$ is the length of the input
string -- each character is examined at most once across the
whitespace-skip, sign-check, and digit-accumulation phases.
\textbf{Space Complexity.} $O(L)$ for the digit-accumulation recursion
stack in the worst case (a string of all digits); this could be reduced
to $O(1)$ by rewriting the accumulation phase as an iterative loop, which
is how production implementations typically do it.
\end{complexitybox}

\begin{takeawaybox}
Not every recursive exercise hides a deep algorithmic insight -- some, like
this one, exist purely to practice expressing a sequential scan (skip,
check, accumulate, stop) as a chain of base-case-and-recurse calls, with
an accumulator threaded through as an explicit parameter rather than a
mutable loop variable. This accumulator-passing style is the same
technique used for tail-recursive functions in functional programming.
\end{takeawaybox}

\section{Reverse a Stack using Recursion}
\label{sec:rec-reverse-stack}

\problemstatement{Given a stack, reverse its order using \textbf{only}
the stack's own \textit{push} and \textit{pop} operations and
\textbf{recursion} -- no second stack, array, or queue is allowed as
auxiliary storage.}

\begin{patternbox}
\emph{``Solve using only recursion, no auxiliary data structure''} is a
direct instruction to use the \textbf{call stack itself} as the only
permitted extra storage: every value temporarily removed from the stack
must be held as a local variable on some recursive call's stack frame,
not in any explicit container. This requires a second, helper recursive
function -- inserting a value back at the \emph{bottom} of a stack is
itself a small recursive sub-problem, since a stack has no native
``insert at the bottom'' operation.
\end{patternbox}

\subsection*{Understanding the problem carefully}

The overall idea: pop the top element off, recursively reverse
\emph{everything else} below it, and then place the popped element back
-- but at the \emph{bottom} of the now-reversed remainder, since
reversal means the original top must end up at the very bottom of the
final stack. Placing a value at the bottom of a stack, using only
\textit{push}/\textit{pop}, is itself solved recursively: pop everything
off one at a time, push the target value once the stack is empty, then
push everything else back on top, in the reverse of the order it was
removed (which recursion handles automatically via the call stack's own
LIFO unwinding).

\begin{ideabox}[Two Cooperating Recursive Functions]
$\textsc{InsertAtBottom}(\textit{stack}, x)$: if $\textit{stack}$ is
empty, push $x$ and return (base case). Otherwise, pop the top into a
local variable $\textit{temp}$, recursively call
$\textsc{InsertAtBottom}(\textit{stack}, x)$ (placing $x$ at the bottom
of what remains), then push $\textit{temp}$ back on top.
$\textsc{ReverseStack}(\textit{stack})$: if $\textit{stack}$ is empty (or
has one element), return (base case). Otherwise, pop the top into
$\textit{temp}$, recursively call
$\textsc{ReverseStack}(\textit{stack})$ (reversing everything below
$\textit{temp}$), then call
$\textsc{InsertAtBottom}(\textit{stack}, \textit{temp})$.
\end{ideabox}

\subsection*{Why local variables on the call stack correctly substitute for an explicit auxiliary structure}

Each recursive call to $\textsc{ReverseStack}$ holds exactly one
popped element in its own local variable $\textit{temp}$, alive for the
duration of that call (including everything that happens inside its
recursive sub-call). Because the call stack itself maintains these
local variables in a strict LIFO order matching the order elements were
popped, and because each element is pushed back (via
$\textsc{InsertAtBottom}$) exactly once as its corresponding call
returns, no information is lost and no second container is ever needed
-- the recursion's own bookkeeping does the job an explicit auxiliary
stack would otherwise have done.

\subsection*{Algorithm}

\begin{enumerate}
  \item $\textsc{InsertAtBottom}(\textit{stack}, x)$:
  \begin{enumerate}
    \item If $\textit{stack}$ is empty: push $x$; return.
    \item Pop top into $\textit{temp}$.
    \item $\textsc{InsertAtBottom}(\textit{stack}, x)$.
    \item Push $\textit{temp}$.
  \end{enumerate}
  \item $\textsc{ReverseStack}(\textit{stack})$:
  \begin{enumerate}
    \item If $\textit{stack}$ is empty: return.
    \item Pop top into $\textit{temp}$.
    \item $\textsc{ReverseStack}(\textit{stack})$.
    \item $\textsc{InsertAtBottom}(\textit{stack}, \textit{temp})$.
  \end{enumerate}
\end{enumerate}

\subsection*{Dry run}

\dryrun{Reverse the stack $[1,2,3]$ (top to bottom: $3,2,1$).}

\begin{itemize}
  \item $\textsc{ReverseStack}$ call 1: pop $3$. Recurse on $[1,2]$.
  \item Call 2: pop $2$. Recurse on $[1]$.
  \item Call 3: pop $1$. Recurse on $[]$ (empty, base case, returns
        immediately).
  \item Back in call 3: $\textsc{InsertAtBottom}([], 1)$: stack is
        empty, push $1$. Stack $=[1]$.
  \item Back in call 2: $\textsc{InsertAtBottom}([1], 2)$: pop $1$ into
        \textit{temp}; recurse $\textsc{InsertAtBottom}([], 2)$: empty,
        push $2$, stack$=[2]$; back in outer call, push \textit{temp}
        ($1$) back: stack$=[2,1]$ (top to bottom: $1,2$).
  \item Back in call 1: $\textsc{InsertAtBottom}([2,1], 3)$: pop $1$
        into \textit{temp}; recurse on $[2]$: pop $2$ into another
        \textit{temp}; recurse on $[]$: push $3$, stack$=[3]$; push
        back the inner \textit{temp} ($2$): stack$=[3,2]$; push back
        the outer \textit{temp} ($1$): stack$=[3,2,1]$.
\end{itemize}

Final stack (top to bottom): $1,2,3$ -- the reverse of the original
$3,2,1$.

\subsection*{C++ implementation}

\begin{lstlisting}
#include <bits/stdc++.h>
using namespace std;

void insertAtBottom(stack<int>& st, int x) {
    if (st.empty()) {
        st.push(x);
        return;
    }
    int temp = st.top();
    st.pop();
    insertAtBottom(st, x);
    st.push(temp);
}

void reverseStack(stack<int>& st) {
    if (st.empty()) return;

    int temp = st.top();
    st.pop();
    reverseStack(st);
    insertAtBottom(st, temp);
}

int main() {
    stack<int> st;
    st.push(1); st.push(2); st.push(3);

    reverseStack(st);

    while (!st.empty()) {
        cout << st.top() << " ";
        st.pop();
    }
    cout << endl;
    return 0;
}
\end{lstlisting}

\begin{complexitybox}
\textbf{Time Complexity.} $\textsc{InsertAtBottom}$ costs $O(k)$ for a
stack of size $k$ (it must pop and re-push every element to reach the
bottom). $\textsc{ReverseStack}$ calls it once for each of the $n$
elements, with the stack shrinking by one each time, giving
\[
  O(1+2+\cdots+n) = O(n^2).
\]
\textbf{Space Complexity.} $O(n)$ for the recursion stack at the deepest
point (both functions' call depths are bounded by $n$).
\end{complexitybox}

\begin{takeawaybox}
``Using only recursion, no auxiliary data structure'' is a direct
instruction to use \emph{local variables on the call stack} as your only
permitted temporary storage. Many such ``stack-only'' problems
(reversing, sorting, this section and the next) need exactly this kind of
helper function that performs an operation a real stack does not natively
support (here, inserting at the bottom), built entirely out of recursive
pop/push/recurse/push-back steps.
\end{takeawaybox}

\section{Sort a Stack using Recursion}
\label{sec:rec-sort-stack}

\problemstatement{Given a stack, sort it in increasing order (so the
smallest element ends up on top) using \textbf{only} \textit{push} and
\textit{pop}, and recursion -- no second stack, array, or queue allowed.}

\begin{patternbox}
This is Section~\ref{sec:rec-reverse-stack}'s ``insert at the bottom''
helper, generalized to ``insert at the \emph{correct sorted position}.''
The overall structure -- pop the top, recursively solve the rest, then
reinsert the popped value using a specialized helper -- is identical;
only the helper's insertion rule changes from ``always go to the
bottom'' to ``go just above everything smaller, and below everything
larger.''
\end{patternbox}

\subsection*{Understanding the problem carefully}

If the remainder of the stack (everything below the current top) is
already correctly sorted, inserting the popped-off top element into its
correct position relative to that sorted remainder finishes the job --
this is exactly insertion sort's core idea, expressed with a stack
playing the role that an array's shifting-elements-right step normally
plays.

\begin{ideabox}[Insert into a Sorted Stack at the Correct Position]
$\textsc{InsertSorted}(\textit{stack}, x)$: if $\textit{stack}$ is empty
\emph{or} its top is $\le x$, push $x$ and return (its correct position
has been found: everything remaining below is $\le x$, satisfying sorted
order with the top of the stack as the largest visible value at this
point... more precisely, the stack is kept increasing from top to
bottom, so $x$ belongs here once it is at least as large as what is
currently on top). Otherwise, pop the top into $\textit{temp}$,
recursively call $\textsc{InsertSorted}(\textit{stack}, x)$ (continue
searching deeper for $x$'s correct position), then push $\textit{temp}$
back on top.
\end{ideabox}

\subsection*{Why this insertion rule keeps the stack sorted at every step}

Assume, inductively, that before inserting $x$, the stack is already
sorted (say, increasing from top to bottom: the smallest element on top).
The helper pops elements off only while they are \emph{larger} than $x$
-- precisely the elements that must end up \emph{above} $x$ once $x$ is
correctly placed. The moment it finds an element $\le x$ (or the stack
empties), that is exactly the right depth to insert $x$, and pushing the
previously-popped larger elements back afterward restores them directly
on top of $x$, in their original relative order (preserved by the
recursive call stack's own LIFO unwinding) -- so the stack remains
sorted after the insertion completes.

\subsection*{Algorithm}

\begin{enumerate}
  \item $\textsc{InsertSorted}(\textit{stack}, x)$:
  \begin{enumerate}
    \item If $\textit{stack}$ is empty or top $\le x$: push $x$; return.
    \item Pop top into $\textit{temp}$.
    \item $\textsc{InsertSorted}(\textit{stack}, x)$.
    \item Push $\textit{temp}$.
  \end{enumerate}
  \item $\textsc{SortStack}(\textit{stack})$:
  \begin{enumerate}
    \item If $\textit{stack}$ is empty: return.
    \item Pop top into $\textit{temp}$.
    \item $\textsc{SortStack}(\textit{stack})$.
    \item $\textsc{InsertSorted}(\textit{stack}, \textit{temp})$.
  \end{enumerate}
\end{enumerate}

\subsection*{Dry run}

\dryrun{Sort the stack $[3,1,2]$ (top to bottom: $2,1,3$).}

\begin{itemize}
  \item Pop $2$. Recurse on $[3,1]$ (top to bottom: $1,3$).
  \item Pop $1$. Recurse on $[3]$.
  \item Pop $3$. Recurse on $[]$ (base case).
  \item $\textsc{InsertSorted}([], 3)$: empty, push. Stack$=[3]$.
  \item Back in the ``pop $1$'' frame: $\textsc{InsertSorted}([3], 1)$:
        top is $3 > 1$, pop $3$ into \textit{temp}; recurse
        $\textsc{InsertSorted}([], 1)$: empty, push $1$, stack$=[1]$;
        push back \textit{temp} ($3$): stack$=[1,3]$ (top to bottom:
        $1,3$ -- sorted).
  \item Back in the ``pop $2$'' frame: $\textsc{InsertSorted}([1,3],
        2)$: top is $1 \le 2$, push $2$ directly. Stack$=[2,1,3]$ (top
        to bottom: $2,1,3$).
\end{itemize}

Final stack, top to bottom: $1,2,3$ -- correctly sorted, smallest on
top.

\subsection*{C++ implementation}

\begin{lstlisting}
#include <bits/stdc++.h>
using namespace std;

void insertSorted(stack<int>& st, int x) {
    if (st.empty() || st.top() <= x) {
        st.push(x);
        return;
    }
    int temp = st.top();
    st.pop();
    insertSorted(st, x);
    st.push(temp);
}

void sortStack(stack<int>& st) {
    if (st.empty()) return;

    int temp = st.top();
    st.pop();
    sortStack(st);
    insertSorted(st, temp);
}

int main() {
    stack<int> st;
    st.push(3); st.push(1); st.push(2);

    sortStack(st);

    while (!st.empty()) {
        cout << st.top() << " ";
        st.pop();
    }
    cout << endl;
    return 0;
}
\end{lstlisting}

\begin{complexitybox}
\textbf{Time Complexity.} $O(n^2)$ in the worst case -- the same
insertion-sort-style bound as Section~\ref{sec:rec-reverse-stack}, since
$\textsc{InsertSorted}$ can pop and re-push up to $O(k)$ elements for a
stack of size $k$, summed over $n$ insertions.
\textbf{Space Complexity.} $O(n)$ for the recursion stack at the deepest
point.
\end{complexitybox}

\begin{takeawaybox}
Once you have solved one ``insert at a specific position using only
recursion'' problem (insert at the bottom, in
Section~\ref{sec:rec-reverse-stack}), solving a closely related one is
just a matter of changing the stopping condition inside the same
recursive shape: stop unconditionally at the bottom for reversal, stop as
soon as a smaller-or-equal element is found for sorting.
\end{takeawaybox}

\section{Power Set: Print All Subsequences}
\label{sec:rec-power-set}

\noindent We now begin the chapter's second family: \textbf{exploration
recursion}, or backtracking. This first problem introduces the single
template -- the \textbf{pick/not-pick} decision tree -- that every other
problem in this group (Sections~\ref{sec:rec-power-set}--\ref{sec:rec-letter-combinations})
specializes or extends.

\problemstatement{Given an array (or string) of $n$ distinct elements,
generate \textbf{every} subsequence (equivalently, every subset, since
order within a subsequence here does not matter for counting purposes)
-- all $2^n$ of them, including the empty subsequence.}

\begin{patternbox}
\emph{``Generate all subsequences/subsets''} is the foundational
backtracking pattern: at every element of the input, there are exactly
\textbf{two} choices -- include it in the current subsequence, or
don't -- and these two choices are completely independent across
different elements. A recursion that branches into exactly these two
calls at every position, one per element, visits every one of the
$2^n$ combinations exactly once.
\end{patternbox}

\subsection*{Understanding the problem carefully}

Think of building a subsequence by walking through the input array left
to right, making one binary decision per element: ``is this element in
my subsequence, or not?'' After making this decision for every one of
the $n$ elements, exactly one specific subsequence has been fully
determined. Since there are $2$ choices at each of $n$ independent
positions, there are exactly $2^n$ distinct sequences of choices --
matching exactly the $2^n$ subsequences a set of $n$ elements has.

\begin{ideabox}[The Pick/Not-Pick Template]
$\textsc{Solve}(\textit{index}, \textit{current})$: if
$\textit{index} = n$ (every element has been decided): record
$\textit{current}$ as one complete subsequence; return (base case).
Otherwise: \textbf{pick} -- append $\textit{arr}[\textit{index}]$ to
$\textit{current}$, recurse with $\textsc{Solve}(\textit{index}+1,
\textit{current})$, then \textbf{undo} the append (remove the last
element, restoring $\textit{current}$ to how it was before this branch).
\textbf{Not-pick} -- recurse with $\textsc{Solve}(\textit{index}+1,
\textit{current})$ unchanged (the element is simply skipped).
\end{ideabox}

\subsection*{Why the explicit ``undo'' step after the pick branch is essential}

If $\textit{current}$ is a single shared mutable list passed by
reference into both branches (rather than a fresh copy made for each),
then after the pick branch's entire recursive exploration finishes and
control returns to this level, $\textit{current}$ must be restored to
exactly the state it had \emph{before} the pick was made -- otherwise the
upcoming not-pick branch would incorrectly start from a list that still
contains the picked element. This pattern -- \textbf{choose}, recurse,
\textbf{un-choose} -- is the defining rhythm of every backtracking
algorithm in this chapter; forgetting the un-choose step is the single
most common bug when implementing backtracking with a shared, mutated
data structure instead of fresh copies.

\subsection*{Algorithm}

\begin{enumerate}
  \item $\textsc{Solve}(\textit{index}, \textit{current})$:
  \begin{enumerate}
    \item If $\textit{index} = n$: add a copy of $\textit{current}$ to
          the result list; return.
    \item Append $\textit{arr}[\textit{index}]$ to $\textit{current}$.
    \item $\textsc{Solve}(\textit{index}+1, \textit{current})$.
    \item Remove the last element of $\textit{current}$ (undo).
    \item $\textsc{Solve}(\textit{index}+1, \textit{current})$
          (not-pick branch, $\textit{current}$ unchanged from before
          step 2).
  \end{enumerate}
  \item Call $\textsc{Solve}(0, [\,])$ to start.
\end{enumerate}

\subsection*{Dry run}

\dryrun{Take $\textit{arr} = [1,2]$ ($n=2$).}

The recursion tree:
\begin{itemize}
  \item $\textsc{Solve}(0, [])$:
  \begin{itemize}
    \item Pick $1$: $\textsc{Solve}(1, [1])$:
    \begin{itemize}
      \item Pick $2$: $\textsc{Solve}(2, [1,2])$: base case, record
            $[1,2]$.
      \item Undo; Not-pick $2$: $\textsc{Solve}(2, [1])$: base case,
            record $[1]$.
    \end{itemize}
    \item Undo; Not-pick $1$: $\textsc{Solve}(1, [])$:
    \begin{itemize}
      \item Pick $2$: $\textsc{Solve}(2, [2])$: base case, record
            $[2]$.
      \item Undo; Not-pick $2$: $\textsc{Solve}(2, [])$: base case,
            record $[]$.
    \end{itemize}
  \end{itemize}
\end{itemize}

Recorded subsequences, in the order generated: $[1,2], [1], [2], []$ --
all $2^2=4$ subsequences of $\{1,2\}$, each exactly once.

\subsection*{C++ implementation}

\begin{lstlisting}
#include <bits/stdc++.h>
using namespace std;

void solve(int index, vector<int>& arr, vector<int>& current,
           vector<vector<int>>& result) {
    if (index == (int)arr.size()) {
        result.push_back(current);
        return;
    }

    current.push_back(arr[index]);      // pick
    solve(index + 1, arr, current, result);
    current.pop_back();                 // undo

    solve(index + 1, arr, current, result); // not-pick
}

vector<vector<int>> powerSet(vector<int>& arr) {
    vector<vector<int>> result;
    vector<int> current;
    solve(0, arr, current, result);
    return result;
}

int main() {
    vector<int> arr = {1, 2};
    for (auto& subset : powerSet(arr)) {
        cout << "[ ";
        for (int x : subset) cout << x << " ";
        cout << "] ";
    }
    cout << endl;
    return 0;
}
\end{lstlisting}

\begin{complexitybox}
\textbf{Time Complexity.} There are $2^n$ leaves in the recursion tree
(one per subsequence), and building each subsequence costs up to $O(n)$
to copy it into the result, giving
\[
  O(n \cdot 2^n).
\]
\textbf{Space Complexity.} $O(n)$ for the recursion depth and the
$\textit{current}$ list at any point in time (excluding the output
storage, which is $O(n \cdot 2^n)$ to hold every generated subsequence).
\end{complexitybox}

\begin{takeawaybox}
The pick/not-pick recursion tree, with an explicit choose/recurse/un-choose
rhythm, is the master template for this entire group of problems.
Combination Sum, Subset Sum, and several other sections ahead are all
this exact same tree, either pruned early (skip branches that can no
longer succeed) or with the choice set widened beyond a simple binary
pick/skip.
\end{takeawaybox}

\section{Check if There Exists a Subsequence with Sum K}
\label{sec:rec-check-subsequence-sum-k}

\problemstatement{Given an array $\textit{arr}$ and an integer $k$,
determine whether \textbf{any} subsequence of $\textit{arr}$ sums to
exactly $k$.}

\begin{patternbox}
This is Section~\ref{sec:rec-power-set}'s pick/not-pick tree, but the
question changes from \emph{``generate every subsequence''} to
\emph{``does at least one satisfy a condition.''} Whenever a backtracking
problem only needs a yes/no answer (rather than every solution), the
recursion should \textbf{stop exploring the instant an answer is found}
-- there is no need to generate the other $2^n-1$ subsequences once one
of them has already confirmed the answer is ``yes.''
\end{patternbox}

\subsection*{Understanding the problem carefully}

Instead of building and recording a list at every leaf, track a running
sum as the pick/not-pick decisions are made, and check it only once
every element has been decided (at the base case). The function should
return a boolean, and -- critically -- the moment either branch reports
\texttt{true}, the calling level should return \texttt{true} immediately
without bothering to explore the other branch at all.

\begin{ideabox}[Same Tree, but Short-Circuit on Success]
$\textsc{Exists}(\textit{index}, \textit{sum})$: if
$\textit{index} = n$: return $(\textit{sum} = k)$. Otherwise: if
$\textsc{Exists}(\textit{index}+1, \textit{sum} +
\textit{arr}[\textit{index}])$ (the pick branch) is \texttt{true}, return
\texttt{true} immediately without checking the not-pick branch at all.
Otherwise, return $\textsc{Exists}(\textit{index}+1, \textit{sum})$ (the
not-pick branch) directly as this call's result.
\end{ideabox}

\subsection*{Why short-circuiting never misses a valid answer}

We only need to know whether \emph{some} subsequence achieves sum $k$ --
not which one, nor how many do. The moment the pick branch reports
success, we already have a definitive ``yes'' for the overall question,
and continuing to explore the not-pick branch (or any other untried
branch) could only ever confirm the same ``yes'' again, never change it
to a ``no.'' Skipping that redundant exploration is therefore always
safe, and in the best case it can save an enormous amount of work,
since a successful subsequence might be found very early in the
traversal while the vast majority of the $2^n$ possibilities are never
even visited.

\subsection*{Algorithm}

\begin{enumerate}
  \item $\textsc{Exists}(\textit{index}, \textit{sum})$:
  \begin{enumerate}
    \item If $\textit{index} = n$: return $(\textit{sum} = k)$.
    \item If $\textsc{Exists}(\textit{index}+1, \textit{sum} +
          \textit{arr}[\textit{index}])$: return \texttt{true}.
    \item Return $\textsc{Exists}(\textit{index}+1, \textit{sum})$.
  \end{enumerate}
  \item Call $\textsc{Exists}(0, 0)$.
\end{enumerate}

\subsection*{Dry run}

\dryrun{Take $\textit{arr} = [1,2,3]$, $k=5$.}

\begin{itemize}
  \item $\textsc{Exists}(0,0)$: pick $1$: $\textsc{Exists}(1,1)$.
  \item $\textsc{Exists}(1,1)$: pick $2$: $\textsc{Exists}(2,3)$.
  \item $\textsc{Exists}(2,3)$: pick $3$: $\textsc{Exists}(3,6)$: base
        case, $6 \ne 5$, return \texttt{false}. Not-pick $3$:
        $\textsc{Exists}(3,3)$: base case, $3 \ne 5$, return
        \texttt{false}. So $\textsc{Exists}(2,3)$ returns
        \texttt{false}.
  \item Back in $\textsc{Exists}(1,1)$: pick branch was \texttt{false},
        try not-pick $2$: $\textsc{Exists}(2,1)$: pick $3$:
        $\textsc{Exists}(3,4)$: base, $4 \ne 5$, false. Not-pick $3$:
        $\textsc{Exists}(3,1)$: base, $1\ne5$, false. So
        $\textsc{Exists}(2,1)$ is false, and $\textsc{Exists}(1,1)$ is
        false.
  \item Back in $\textsc{Exists}(0,0)$: pick branch was \texttt{false},
        try not-pick $1$: $\textsc{Exists}(1,0)$: pick $2$:
        $\textsc{Exists}(2,2)$: pick $3$: $\textsc{Exists}(3,5)$: base,
        $5=5$! Return \texttt{true}. This propagates straight back up:
        $\textsc{Exists}(2,2)$ returns \texttt{true} immediately (no
        need to check not-pick $3$); $\textsc{Exists}(1,0)$ returns
        \texttt{true} immediately; $\textsc{Exists}(0,0)$ returns
        \texttt{true}.
\end{itemize}

The answer is \texttt{true}: the subsequence $\{2,3\}$ sums to $5$, and
the recursion stopped exploring further branches the instant it was
found.

\subsection*{C++ implementation}

\begin{lstlisting}
#include <bits/stdc++.h>
using namespace std;

bool exists(int index, int sum, vector<int>& arr, int k) {
    if (index == (int)arr.size()) return sum == k;

    if (exists(index + 1, sum + arr[index], arr, k)) return true;
    return exists(index + 1, sum, arr, k);
}

int main() {
    vector<int> arr = {1, 2, 3};
    cout << (exists(0, 0, arr, 5) ? "true" : "false") << endl;
    return 0;
}
\end{lstlisting}

\begin{complexitybox}
\textbf{Time Complexity.} $O(2^n)$ in the worst case (no qualifying
subsequence exists, so every branch must be explored), but often much
faster in practice due to short-circuiting on the first success.
\textbf{Space Complexity.} $O(n)$ for the recursion stack.
\end{complexitybox}

\begin{takeawaybox}
``Does there exist'' questions should always short-circuit the moment a
valid answer is found -- propagate a \texttt{true} result immediately up
through every level of the recursion, skipping any sibling branch that
has not yet been explored. This is a small but important variation on
the pick/not-pick template: the same tree shape, but with an early-exit
discipline that the ``generate everything'' version of
Section~\ref{sec:rec-power-set} does not need.
\end{takeawaybox}

\section{Count Subsequences with Sum K}
\label{sec:rec-count-subsequence-sum-k}

\problemstatement{Given an array $\textit{arr}$ and an integer $k$, count
\textbf{how many} subsequences of $\textit{arr}$ sum to exactly $k$.}

\begin{patternbox}
The third variation on Section~\ref{sec:rec-power-set}'s pick/not-pick
tree: from \emph{``generate all''} (Section~\ref{sec:rec-power-set}) to
\emph{``does one exist''} (Section~\ref{sec:rec-check-subsequence-sum-k})
to, now, \emph{``count how many.''} Counting needs neither full
generation nor early short-circuiting -- it needs both branches'
\textbf{counts} added together.
\end{patternbox}

\subsection*{Understanding the problem carefully}

Every complete pick/not-pick decision sequence produces exactly one
subsequence, and the base case can directly check whether \emph{that}
subsequence's sum equals $k$, contributing $1$ to the count if so, $0$
otherwise. Since the pick and not-pick branches partition all
subsequences into two disjoint groups (those that include
$\textit{arr}[\textit{index}]$ and those that don't), the total count is
simply the \textbf{sum} of the counts returned by both branches -- no
branch can be skipped, since both might contain qualifying subsequences.

\begin{ideabox}[Sum the Counts from Both Branches]
$\textsc{CountSum}(\textit{index}, \textit{sum})$: if
$\textit{index} = n$: return $1$ if $\textit{sum} = k$, else $0$.
Otherwise: return $\textsc{CountSum}(\textit{index}+1, \textit{sum} +
\textit{arr}[\textit{index}]) + \textsc{CountSum}(\textit{index}+1,
\textit{sum})$.
\end{ideabox}

\subsection*{Why summing (not short-circuiting) is required here}

Unlike Section~\ref{sec:rec-check-subsequence-sum-k}, where finding one
success already fully answers the yes/no question, here a successful
pick branch tells us nothing about how many \emph{additional} successes
might be hiding in the not-pick branch -- both must be fully explored and
their individual counts combined, since the two branches' qualifying
subsequences are entirely different (disjoint) subsequences, both of
which need to be counted toward the final total.

\subsection*{Algorithm}

\begin{enumerate}
  \item $\textsc{CountSum}(\textit{index}, \textit{sum})$:
  \begin{enumerate}
    \item If $\textit{index} = n$: return $1$ if $\textit{sum}=k$, else
          $0$.
    \item Return $\textsc{CountSum}(\textit{index}+1, \textit{sum} +
          \textit{arr}[\textit{index}]) +
          \textsc{CountSum}(\textit{index}+1, \textit{sum})$.
  \end{enumerate}
  \item Call $\textsc{CountSum}(0,0)$.
\end{enumerate}

\subsection*{Dry run}

\dryrun{Take $\textit{arr} = [1,2,1]$, $k=2$.}

\begin{itemize}
  \item $\textsc{CountSum}(0,0)$: pick $1$:
        $\textsc{CountSum}(1,1)$; not-pick: $\textsc{CountSum}(1,0)$.
  \item $\textsc{CountSum}(1,1)$ (came from picking index $0$): pick
        $2$: $\textsc{CountSum}(2,3)$; not-pick:
        $\textsc{CountSum}(2,1)$.
  \begin{itemize}
    \item $\textsc{CountSum}(2,3)$: pick $1$: base, sum$=4\ne2$,
          return $0$. Not-pick: base, sum$=3\ne2$, return $0$. Total
          $0$.
    \item $\textsc{CountSum}(2,1)$: pick $1$: base, sum$=2=2$! return
          $1$. Not-pick: base, sum$=1\ne2$, return $0$. Total $1$.
  \end{itemize}
  \item $\textsc{CountSum}(1,1)=0+1=1$.
  \item $\textsc{CountSum}(1,0)$ (came from not-picking index $0$):
        pick $2$: $\textsc{CountSum}(2,2)$; not-pick:
        $\textsc{CountSum}(2,0)$.
  \begin{itemize}
    \item $\textsc{CountSum}(2,2)$: pick $1$: base, sum$=3\ne2$,
          return $0$. Not-pick: base, sum$=2=2$! return $1$. Total $1$.
    \item $\textsc{CountSum}(2,0)$: pick $1$: base, sum$=1\ne2$,
          return $0$. Not-pick: base, sum$=0\ne2$, return $0$. Total
          $0$.
  \end{itemize}
  \item $\textsc{CountSum}(1,0)=1+0=1$.
  \item $\textsc{CountSum}(0,0)=1+1=2$.
\end{itemize}

The answer is $2$: the qualifying subsequences are $\{1,1\}$ (indices
$0,2$) and $\{2\}$ (index $1$ alone).

\subsection*{C++ implementation}

\begin{lstlisting}
#include <bits/stdc++.h>
using namespace std;

int countSum(int index, int sum, vector<int>& arr, int k) {
    if (index == (int)arr.size()) return sum == k ? 1 : 0;

    int pick = countSum(index + 1, sum + arr[index], arr, k);
    int notPick = countSum(index + 1, sum, arr, k);
    return pick + notPick;
}

int main() {
    vector<int> arr = {1, 2, 1};
    cout << "Count: " << countSum(0, 0, arr, 2) << endl;
    return 0;
}
\end{lstlisting}

\begin{complexitybox}
\textbf{Time Complexity.} $O(2^n)$ -- every one of the $2^n$
subsequences must be considered, since any of them might contribute to
the count.
\textbf{Space Complexity.} $O(n)$ for the recursion stack.
\end{complexitybox}

\begin{takeawaybox}
The three variations seen so far on the same pick/not-pick tree --
generate everything (Section~\ref{sec:rec-power-set}), short-circuit on
the first success (Section~\ref{sec:rec-check-subsequence-sum-k}), and
sum counts from both branches (this section) -- are worth holding side by
side as a single mental checklist: \emph{generate, exists, or count} is
usually the very first distinction to make about any new backtracking
problem, since it determines how the two branches' results should be
combined at every level.
\end{takeawaybox}

\section{Subset Sums I (Print All Subset Sums)}
\label{sec:rec-subset-sum-1}

\problemstatement{Given an array $\textit{arr}$ of $n$ integers, compute
the \textbf{sum} of every one of its $2^n$ subsets, and return the list
of all such sums (typically in sorted order).}

\begin{patternbox}
A direct simplification of Section~\ref{sec:rec-power-set}'s
pick/not-pick tree: instead of recording the entire subsequence at each
leaf, we only need to record its \textbf{sum} -- a single number threaded
through the recursion as a running total, rather than a list that must
be built and torn down at every level.
\end{patternbox}

\subsection*{Understanding the problem carefully}

Exactly as in Section~\ref{sec:rec-power-set}, every element independently
contributes a pick/not-pick decision. The only change is what we choose
to track: instead of appending the picked element to a growing list (and
later removing it on backtrack), we simply add its value to a running
sum parameter (and, since this is passed by value rather than by
reference, there is nothing to explicitly undo -- each recursive call
gets its own independent copy of the sum).

\begin{ideabox}[Thread a Running Sum Instead of a List]
$\textsc{Solve}(\textit{index}, \textit{sum})$: if
$\textit{index} = n$: record $\textit{sum}$ as one subset's total;
return. Otherwise: $\textsc{Solve}(\textit{index}+1, \textit{sum} +
\textit{arr}[\textit{index}])$ (pick), then
$\textsc{Solve}(\textit{index}+1, \textit{sum})$ (not-pick).
\end{ideabox}

\subsection*{Why no explicit undo step is needed here}

In Section~\ref{sec:rec-power-set}, $\textit{current}$ was a single
shared, mutated list, so the pick branch's append had to be explicitly
undone before exploring the not-pick branch. Here, $\textit{sum}$ is
passed \textbf{by value} (a plain number, copied into each recursive
call rather than referenced and mutated in place), so the pick branch's
modified sum simply ceases to exist once that branch's call returns --
the not-pick branch automatically continues with the original,
unmodified sum from this level, with no explicit ``undo'' arithmetic
required. This is a useful general lesson: whether backtracking needs an
explicit undo step depends entirely on whether the state being threaded
through is mutated in place (needing undo) or passed by value
(automatically restored).

\subsection*{Algorithm}

\begin{enumerate}
  \item $\textsc{Solve}(\textit{index}, \textit{sum})$:
  \begin{enumerate}
    \item If $\textit{index}=n$: append $\textit{sum}$ to the result
          list; return.
    \item $\textsc{Solve}(\textit{index}+1, \textit{sum} +
          \textit{arr}[\textit{index}])$.
    \item $\textsc{Solve}(\textit{index}+1, \textit{sum})$.
  \end{enumerate}
  \item Call $\textsc{Solve}(0,0)$; sort the result list before
        returning (if sorted output is required).
\end{enumerate}

\subsection*{Dry run}

\dryrun{Take $\textit{arr} = [2,3]$.}

\begin{itemize}
  \item $\textsc{Solve}(0,0)$: pick $2$: $\textsc{Solve}(1,2)$.
  \begin{itemize}
    \item Pick $3$: $\textsc{Solve}(2,5)$: base, record $5$.
    \item Not-pick $3$: $\textsc{Solve}(2,2)$: base, record $2$.
  \end{itemize}
  \item Not-pick $2$: $\textsc{Solve}(1,0)$.
  \begin{itemize}
    \item Pick $3$: $\textsc{Solve}(2,3)$: base, record $3$.
    \item Not-pick $3$: $\textsc{Solve}(2,0)$: base, record $0$.
  \end{itemize}
\end{itemize}

Recorded sums: $5,2,3,0$; sorted: $0,2,3,5$ -- matching the four subsets
$\{\}, \{2\}, \{3\}, \{2,3\}$.

\subsection*{C++ implementation}

\begin{lstlisting}
#include <bits/stdc++.h>
using namespace std;

void solve(int index, int sum, vector<int>& arr, vector<int>& result) {
    if (index == (int)arr.size()) {
        result.push_back(sum);
        return;
    }
    solve(index + 1, sum + arr[index], arr, result);
    solve(index + 1, sum, arr, result);
}

vector<int> subsetSums(vector<int>& arr) {
    vector<int> result;
    solve(0, 0, arr, result);
    sort(result.begin(), result.end());
    return result;
}

int main() {
    vector<int> arr = {2, 3};
    for (int s : subsetSums(arr)) cout << s << " ";
    cout << endl;
    return 0;
}
\end{lstlisting}

\begin{complexitybox}
\textbf{Time Complexity.} $O(2^n)$ to generate all sums, plus
$O(2^n \log(2^n)) = O(n \cdot 2^n)$ if sorting the output is required.
\textbf{Space Complexity.} $O(n)$ for the recursion stack, plus
$O(2^n)$ for the output.
\end{complexitybox}

\begin{takeawaybox}
Whether a backtracking branch needs an explicit ``undo'' depends on
whether its state is mutated in place (a shared list, needing an undo)
or passed by value (a plain number or immutable structure, automatically
restored when the branch returns). Recognizing which kind of state a
problem needs is often the deciding factor in how cleanly the recursion
can be written.
\end{takeawaybox}

\section{Subset Sums II (Unique Subsets with Duplicates Present)}
\label{sec:rec-subset-sum-2}

\problemstatement{Given an array $\textit{arr}$ of $n$ integers that
\textbf{may contain duplicate values}, return every distinct subset
(the power set), with \textbf{no duplicate subsets} in the output, even
though duplicate elements exist in the input.}

\begin{patternbox}
The keyword \emph{``may contain duplicates,''} combined with
\emph{``no duplicate subsets in the output,''} is the trigger for this
chapter's other major backtracking template, distinct from pick/not-pick:
\textbf{sort first, then iterate through choices at each recursive level,
explicitly skipping any choice equal to the immediately previous choice
already tried at the \emph{same} level.} This ``sort, then skip
same-level duplicates'' technique is reused in every later problem that
needs to handle duplicate input values
(Sections~\ref{sec:rec-combination-sum-2}--\ref{sec:rec-combination-sum-3}).
\end{patternbox}

\subsection*{Understanding the problem carefully}

If the input has duplicate values, the plain pick/not-pick tree from
Section~\ref{sec:rec-power-set} would generate the \emph{same} subset
multiple times -- for example, with $\textit{arr}=[1,2,2]$, picking
``the first $2$'' and picking ``the second $2$'' (while skipping the
other) produce two different leaves in the recursion tree, but the
exact same subset $\{1,2\}$. To avoid this, we change the recursion's
shape: instead of a binary pick/not-pick at every fixed index, we
\textbf{sort} the array first (so equal values become adjacent), and at
each recursive level, we loop through the remaining choices, but skip
any choice that is equal to the choice immediately before it \emph{at
this same level} (not across different levels, where repeats are
exactly what we want, to build subsets containing multiple copies of a
repeated value).

\begin{ideabox}[Sort First, Then Skip Same-Level Duplicates]
Sort $\textit{arr}$. $\textsc{Solve}(\textit{start}, \textit{current})$:
record a copy of $\textit{current}$ as one valid subset immediately
(every prefix explored is itself a valid, distinct subset, not just the
leaves). For $i$ from $\textit{start}$ to $n-1$: if $i > \textit{start}$
and $\textit{arr}[i] = \textit{arr}[i-1]$, skip this $i$ (this exact
value has already been tried as the \emph{first} new element added at
this level, so trying it again here would only reproduce a subset
already generated). Otherwise: append $\textit{arr}[i]$ to
$\textit{current}$; recurse $\textsc{Solve}(i+1, \textit{current})$;
remove $\textit{arr}[i]$ from $\textit{current}$ (undo).
\end{ideabox}

\subsection*{Why skipping only at the same level (not across levels) is correct}

The key distinction: $\textit{arr}[i] = \textit{arr}[i-1]$ being skipped
only applies when both are being considered as candidates \emph{to fill
the same ``next new element'' slot} within the same call to
$\textsc{Solve}$ -- i.e.\ within the same \texttt{for} loop, at the same
recursion depth. A repeated value is never forbidden from appearing
\emph{multiple times within one subset} (e.g.\ $\{2,2\}$ from
$\textit{arr}=[1,2,2]$ is a perfectly valid, distinct subset) -- what is
forbidden is generating the \emph{same subset twice} by choosing ``the
first $2$ at this position'' in one branch and ``the second $2$ at this
same position'' in a sibling branch, when both choices would lead to
identical subsequent subsets. Since the array is sorted, any two equal
values are guaranteed to be adjacent in the loop, so checking
$\textit{arr}[i] = \textit{arr}[i-1]$ exactly identifies this situation.

\subsection*{Algorithm}

\begin{enumerate}
  \item Sort $\textit{arr}$.
  \item $\textsc{Solve}(\textit{start}, \textit{current})$:
  \begin{enumerate}
    \item Append a copy of $\textit{current}$ to the result.
    \item For $i$ from $\textit{start}$ to $n-1$:
    \begin{itemize}
      \item If $i>\textit{start}$ and $\textit{arr}[i] =
            \textit{arr}[i-1]$: skip (continue to next $i$).
      \item Append $\textit{arr}[i]$ to $\textit{current}$; recurse
            $\textsc{Solve}(i+1, \textit{current})$; remove
            $\textit{arr}[i]$ (undo).
    \end{itemize}
  \end{enumerate}
  \item Call $\textsc{Solve}(0, [\,])$.
\end{enumerate}

\subsection*{Dry run}

\dryrun{Take $\textit{arr} = [1,2,2]$ (already sorted).}

\begin{itemize}
  \item $\textsc{Solve}(0,[])$: record $[]$. Loop $i=0$ ($1$): append.
        $\textsc{Solve}(1,[1])$. Undo.
  \begin{itemize}
    \item $\textsc{Solve}(1,[1])$: record $[1]$. Loop $i=1$ ($2$):
          append. $\textsc{Solve}(2,[1,2])$. Undo. $i=2$ ($2$): $i>
          \textit{start}=1$ and $\textit{arr}[2]=\textit{arr}[1]=2$:
          skip.
    \begin{itemize}
      \item $\textsc{Solve}(2,[1,2])$: record $[1,2]$. Loop $i=2$
            ($2$): append. $\textsc{Solve}(3,[1,2,2])$: record
            $[1,2,2]$. Undo.
    \end{itemize}
  \end{itemize}
  \item Loop continues at top level: $i=1$ ($2$): $i>\textit{start}=0$,
        but $\textit{arr}[1] \ne \textit{arr}[0]$ ($2\ne1$), no skip.
        Append. $\textsc{Solve}(2,[2])$. Undo.
  \begin{itemize}
    \item $\textsc{Solve}(2,[2])$: record $[2]$. Loop $i=2$ ($2$):
          append. $\textsc{Solve}(3,[2,2])$: record $[2,2]$. Undo.
  \end{itemize}
  \item $i=2$ at top level: $i>\textit{start}=0$ and
        $\textit{arr}[2]=\textit{arr}[1]=2$: skip.
\end{itemize}

Recorded subsets: $[], [1], [1,2], [1,2,2], [2], [2,2]$ -- exactly the
$6$ distinct subsets of $\{1,2,2\}$ (note there are $6$, not $2^3=8$,
because $\{2\}$ and $\{2,2\}$ each appear in only one form despite the
duplicate $2$).

\subsection*{C++ implementation}

\begin{lstlisting}
#include <bits/stdc++.h>
using namespace std;

void solve(int start, vector<int>& arr, vector<int>& current,
           vector<vector<int>>& result) {
    result.push_back(current);

    for (int i = start; i < (int)arr.size(); i++) {
        if (i > start && arr[i] == arr[i - 1]) continue; // skip duplicate

        current.push_back(arr[i]);
        solve(i + 1, arr, current, result);
        current.pop_back();
    }
}

vector<vector<int>> subsetsWithDup(vector<int>& arr) {
    sort(arr.begin(), arr.end());
    vector<vector<int>> result;
    vector<int> current;
    solve(0, arr, current, result);
    return result;
}

int main() {
    vector<int> arr = {1, 2, 2};
    for (auto& subset : subsetsWithDup(arr)) {
        cout << "[ ";
        for (int x : subset) cout << x << " ";
        cout << "] ";
    }
    cout << endl;
    return 0;
}
\end{lstlisting}

\begin{complexitybox}
\textbf{Time Complexity.} Sorting costs $O(n \log n)$. The recursion
visits each distinct subset once, and building/copying each costs up to
$O(n)$, giving an overall time complexity bounded by
\[
  O(n \cdot 2^n)
\]
in the worst case (no duplicates, so all $2^n$ subsets are distinct).
\textbf{Space Complexity.} $O(n)$ for the recursion depth and
$\textit{current}$ list.
\end{complexitybox}

\begin{takeawaybox}
``Sort first, then skip a choice equal to the immediately previous
choice tried at the same recursion level'' is the standard fix for
duplicate-avoidance in backtracking, and it generalizes far beyond
subsets: the same technique, applied to the choice loop inside
Combination Sum II (Section~\ref{sec:rec-combination-sum-2}) and
Combination Sum III (Section~\ref{sec:rec-combination-sum-3}), prevents
duplicate combinations there too.
\end{takeawaybox}

\section{Combination Sum}
\label{sec:rec-combination-sum}

\problemstatement{Given an array $\textit{candidates}$ of \textbf{distinct}
positive integers and a target $\textit{target}$, return every unique
combination of candidates that sums to $\textit{target}$. The
\textbf{same candidate may be used an unlimited number of times} within
one combination.}

\begin{patternbox}
\emph{``The same element may be reused unlimited times''} is the
keyword that distinguishes this from Section~\ref{sec:rec-power-set}'s
plain pick/not-pick tree: the \textbf{pick} branch must \textbf{stay at
the same index} (since this element remains available for reuse) rather
than advancing to the next index, while the \textbf{not-pick} branch
still advances (moving permanently past this candidate, never to
reconsider it).
\end{patternbox}

\subsection*{Understanding the problem carefully}

At each index, we still face a binary choice -- use this candidate (and
possibly use it again later) or move past it forever -- but ``use this
candidate'' no longer means ``use it exactly once and move on'': because
reuse is unlimited, choosing to include $\textit{candidates}[\textit{index}]$
should recurse back into the \emph{same} index, leaving open the
possibility of including it again. The recursion only definitively moves
past an index once the not-pick branch is taken.

\begin{ideabox}[Pick Stays, Not-Pick Advances]
$\textsc{Solve}(\textit{index}, \textit{remaining}, \textit{current})$:
if $\textit{remaining}=0$: record $\textit{current}$; return. If
$\textit{index}=n$ or $\textit{remaining}<0$: return (failed branch,
either ran out of candidates or overshot the target). Otherwise: append
$\textit{candidates}[\textit{index}]$ to $\textit{current}$; recurse
$\textsc{Solve}(\textit{index}, \textit{remaining} -
\textit{candidates}[\textit{index}], \textit{current})$ (pick, staying
at the same index); remove the appended element (undo); recurse
$\textsc{Solve}(\textit{index}+1, \textit{remaining}, \textit{current})$
(not-pick, advancing).
\end{ideabox}

\subsection*{Why staying at the same index correctly models unlimited reuse}

If the pick branch advanced to $\textit{index}+1$ as in
Section~\ref{sec:rec-power-set}, each candidate could be used at most
once per combination -- exactly the wrong behavior for this problem.
Staying at $\textit{index}$ after picking means the very next recursive
call still has $\textit{candidates}[\textit{index}]$ available to pick
\emph{again}, and can keep doing so for as many repetitions as the
remaining target allows, only ever moving past index $\textit{index}$
once the not-pick branch is explicitly taken (which represents
``permanently done with this candidate, for this branch of the search'').
The $\textit{remaining}<0$ check ensures the recursion does not loop
forever picking a positive candidate repeatedly past the point where it
could still sum to the target -- it is the base-case escape that makes
termination certain despite the same index being revisited arbitrarily
many times.

\subsection*{Algorithm}

\begin{enumerate}
  \item $\textsc{Solve}(\textit{index}, \textit{remaining}, \textit{current})$:
  \begin{enumerate}
    \item If $\textit{remaining}=0$: record $\textit{current}$; return.
    \item If $\textit{index}=n$ or $\textit{remaining}<0$: return.
    \item Append $\textit{candidates}[\textit{index}]$;
          $\textsc{Solve}(\textit{index}, \textit{remaining} -
          \textit{candidates}[\textit{index}], \textit{current})$;
          remove (undo).
    \item $\textsc{Solve}(\textit{index}+1, \textit{remaining},
          \textit{current})$.
  \end{enumerate}
  \item Call $\textsc{Solve}(0, \textit{target}, [\,])$.
\end{enumerate}

\subsection*{Dry run}

\dryrun{Take $\textit{candidates} = [2,3]$, $\textit{target}=5$.}

\begin{itemize}
  \item $\textsc{Solve}(0,5,[])$: pick $2$:
        $\textsc{Solve}(0,3,[2])$ (stays at index $0$).
  \begin{itemize}
    \item Pick $2$ again: $\textsc{Solve}(0,1,[2,2])$. Pick $2$ again:
          $\textsc{Solve}(0,-1,[2,2,2])$: $\textit{remaining}<0$,
          return (fail). Undo. Not-pick: $\textsc{Solve}(1,1,[2,2])$:
          pick $3$: $\textsc{Solve}(1,-2,\ldots)$: fail. Undo.
          Not-pick: $\textsc{Solve}(2,1,[2,2])$: $\textit{index}=n=2$,
          return (fail).
    \item Undo back to $[2]$. Not-pick $3$ at this level:
          $\textsc{Solve}(1,3,[2])$: pick $3$:
          $\textsc{Solve}(1,0,[2,3])$: $\textit{remaining}=0$! Record
          $[2,3]$.
  \end{itemize}
  \item Back at top: undo to $[]$. Not-pick $2$:
        $\textsc{Solve}(1,5,[])$: pick $3$: $\textsc{Solve}(1,2,[3])$:
        pick $3$ again: $\textsc{Solve}(1,-1,[3,3])$: fail. Not-pick:
        $\textsc{Solve}(2,2,[3])$: $\textit{index}=n$, fail.
\end{itemize}

The only recorded combination is $[2,3]$. (Note: in a full trace,
$[2,2,\ldots]$ paths never reach exactly $0$ remaining and so never
record anything, which the dry run above confirms.)

\subsection*{C++ implementation}

\begin{lstlisting}
#include <bits/stdc++.h>
using namespace std;

void solve(int index, int remaining, vector<int>& candidates,
           vector<int>& current, vector<vector<int>>& result) {
    if (remaining == 0) {
        result.push_back(current);
        return;
    }
    if (index == (int)candidates.size() || remaining < 0) return;

    current.push_back(candidates[index]);
    solve(index, remaining - candidates[index], candidates, current, result);
    current.pop_back();

    solve(index + 1, remaining, candidates, current, result);
}

vector<vector<int>> combinationSum(vector<int>& candidates, int target) {
    vector<vector<int>> result;
    vector<int> current;
    solve(0, target, candidates, current, result);
    return result;
}

int main() {
    vector<int> candidates = {2, 3, 6, 7};
    for (auto& comb : combinationSum(candidates, 7)) {
        cout << "[ ";
        for (int x : comb) cout << x << " ";
        cout << "] ";
    }
    cout << endl;
    return 0;
}
\end{lstlisting}

\begin{complexitybox}
\textbf{Time Complexity.} Difficult to bound tightly in general (it
depends on the target and candidate values), but is exponential in the
worst case, roughly $O(2^{\textit{target}})$ for the branching factor at
each remaining-target unit; in practice it is bounded by the number of
valid combinations times their average length.
\textbf{Space Complexity.} $O(\textit{target}/\min(\textit{candidates}))$
for the maximum recursion depth (the most times the smallest candidate
could be picked before exceeding the target).
\end{complexitybox}

\begin{takeawaybox}
``Unlimited reuse of the same element'' converts the pick branch's
target index from ``advance to $\textit{index}+1$'' to ``stay at
$\textit{index}$.'' This single change is the entire difference between
Combination Sum (this section) and the no-reuse subset/combination
problems before and after it -- worth recognizing on sight whenever a
problem statement explicitly allows repeated use of the same value.
\end{takeawaybox}

\section{Combination Sum II}
\label{sec:rec-combination-sum-2}

\problemstatement{Given an array $\textit{candidates}$ that \textbf{may
contain duplicate values} and a target, return every unique combination
that sums to the target, where \textbf{each individual element of the
array may be used at most once} (no reuse), and the output must contain
no duplicate combinations.}

\begin{patternbox}
This problem directly combines two techniques already derived in this
chapter: Section~\ref{sec:rec-subset-sum-2}'s ``sort, then skip a choice
equal to the previous choice at the same level'' duplicate-avoidance, and
Section~\ref{sec:rec-combination-sum}'s general combination-building
recursion -- but \textbf{without} that section's ``stay at the same
index'' reuse rule, since each element here may only be used once.
\end{patternbox}

\subsection*{Understanding the problem carefully}

Because reuse is forbidden, the recursion always advances past the
current index after picking it (exactly like
Section~\ref{sec:rec-subset-sum-2}, never like
Section~\ref{sec:rec-combination-sum}). Because duplicate values are
present, we sort first and skip same-level duplicate choices, exactly as
in Section~\ref{sec:rec-subset-sum-2}, to avoid generating the same
combination more than once. One further optimization becomes available
specifically because we additionally have a numeric target and a sorted
array: the moment a candidate's value alone exceeds the remaining
target, every later candidate (being at least as large, since the array
is sorted) must also exceed it, so the entire rest of the loop at this
level can be abandoned immediately, rather than checked one by one.

\begin{ideabox}[Sort, Skip Same-Level Duplicates, and Break Early]
Sort $\textit{candidates}$. $\textsc{Solve}(\textit{start},
\textit{remaining}, \textit{current})$: if $\textit{remaining}=0$:
record $\textit{current}$; return. For $i$ from $\textit{start}$ to
$n-1$: if $i>\textit{start}$ and $\textit{candidates}[i] =
\textit{candidates}[i-1]$: skip (same-level duplicate). If
$\textit{candidates}[i] > \textit{remaining}$: \textbf{break} out of the
loop entirely (pruning -- no further $i$ in this sorted loop can help).
Otherwise: append $\textit{candidates}[i]$; recurse
$\textsc{Solve}(i+1, \textit{remaining} - \textit{candidates}[i],
\textit{current})$; remove (undo).
\end{ideabox}

\subsection*{Why breaking (not just skipping) on an over-large candidate is valid}

This is a direct consequence of sortedness, not a heuristic: since
$\textit{candidates}$ is sorted in increasing order, once
$\textit{candidates}[i] > \textit{remaining}$, every subsequent
$\textit{candidates}[j]$ for $j>i$ satisfies $\textit{candidates}[j] \ge
\textit{candidates}[i] > \textit{remaining}$ too -- so none of them could
possibly lead to a valid combination either, at this level. Breaking the
loop entirely (rather than merely skipping this one $i$ and checking the
next) is therefore a safe pruning step that discards an entire suffix of
guaranteed-failing branches in $O(1)$, rather than wastefully visiting
and individually rejecting each one.

\subsection*{Algorithm}

\begin{enumerate}
  \item Sort $\textit{candidates}$.
  \item $\textsc{Solve}(\textit{start}, \textit{remaining}, \textit{current})$:
  \begin{enumerate}
    \item If $\textit{remaining}=0$: record $\textit{current}$; return.
    \item For $i$ from $\textit{start}$ to $n-1$:
    \begin{itemize}
      \item If $i>\textit{start}$ and $\textit{candidates}[i] =
            \textit{candidates}[i-1]$: continue (skip).
      \item If $\textit{candidates}[i] > \textit{remaining}$: break.
      \item Append $\textit{candidates}[i]$; recurse
            $\textsc{Solve}(i+1, \textit{remaining} -
            \textit{candidates}[i], \textit{current})$; remove (undo).
    \end{itemize}
  \end{enumerate}
  \item Call $\textsc{Solve}(0, \textit{target}, [\,])$.
\end{enumerate}

\subsection*{Dry run}

\dryrun{Take $\textit{candidates} = [1,1,2]$ (sorted),
$\textit{target}=3$.}

\begin{itemize}
  \item $\textsc{Solve}(0,3,[])$: $i=0$ ($1$): append, recurse
        $\textsc{Solve}(1,2,[1])$. Undo.
  \begin{itemize}
    \item $\textsc{Solve}(1,2,[1])$: $i=1$ ($1$): $i=\textit{start}$,
          no skip check needed. Append, recurse
          $\textsc{Solve}(2,1,[1,1])$. Undo.
    \begin{itemize}
      \item $\textsc{Solve}(2,1,[1,1])$: $i=2$ ($2$): $2 > 1$,
            \textbf{break} immediately.
    \end{itemize}
    \item $i=2$ ($2$) at this level: append, recurse
          $\textsc{Solve}(3,0,[1,2])$: $\textit{remaining}=0$! Record
          $[1,2]$.
  \end{itemize}
  \item Back at top level: $i=1$ ($1$): $i>\textit{start}=0$ and
        $\textit{candidates}[1]=\textit{candidates}[0]$: skip.
  \item $i=2$ ($2$): append, recurse $\textsc{Solve}(3,1,[2])$:
        $\textit{remaining}=1 \ne 0$, and the loop from $i=3$ to $n-1$
        is empty (no candidates left); this branch records nothing.
\end{itemize}

Recorded combinations: $[1,1,1]$? -- wait, recheck: we recorded $[1,2]$
once. The full correct output for this classic instance is
$\{[1,1,1]\text{-like attempts that failed}, [1,2]\}$; tracing carefully
as above, the only successful combination found is $[1,2]$, matching the
expected unique result for target $3$ from $\{1,1,2\}$ (the combination
$1+1+2=4 \ne 3$ does not qualify; $1+2=3$ does, exactly once despite two
copies of $1$ being available).

\subsection*{C++ implementation}

\begin{lstlisting}
#include <bits/stdc++.h>
using namespace std;

void solve(int start, int remaining, vector<int>& candidates,
           vector<int>& current, vector<vector<int>>& result) {
    if (remaining == 0) {
        result.push_back(current);
        return;
    }
    for (int i = start; i < (int)candidates.size(); i++) {
        if (i > start && candidates[i] == candidates[i - 1]) continue;
        if (candidates[i] > remaining) break;

        current.push_back(candidates[i]);
        solve(i + 1, remaining - candidates[i], candidates, current, result);
        current.pop_back();
    }
}

vector<vector<int>> combinationSum2(vector<int>& candidates, int target) {
    sort(candidates.begin(), candidates.end());
    vector<vector<int>> result;
    vector<int> current;
    solve(0, target, candidates, current, result);
    return result;
}

int main() {
    vector<int> candidates = {1, 1, 2};
    for (auto& comb : combinationSum2(candidates, 3)) {
        cout << "[ ";
        for (int x : comb) cout << x << " ";
        cout << "] ";
    }
    cout << endl;
    return 0;
}
\end{lstlisting}

\begin{complexitybox}
\textbf{Time Complexity.} $O(2^n)$ in the worst case (every subset
potentially checked), though sorting plus the early-break pruning
typically eliminates a substantial fraction of branches in practice; an
additional $O(n \log n)$ for the initial sort.
\textbf{Space Complexity.} $O(n)$ for the recursion depth.
\end{complexitybox}

\begin{takeawaybox}
Combination Sum II is a direct composition of two already-derived
techniques -- no-reuse advancing recursion plus sort-and-skip duplicate
avoidance -- with one additional sorted-array pruning trick (breaking
early once a candidate exceeds the remaining target) layered on top.
Recognizing a new problem as ``technique A combined with technique B''
is, again, the fastest path to a correct solution.
\end{takeawaybox}

\section{Combination Sum III}
\label{sec:rec-combination-sum-3}

\problemstatement{Find all valid combinations of exactly $k$
\textbf{distinct} numbers, chosen from $1$ through $9$ (each usable at
most once across all combinations, i.e.\ no repeats within one
combination), that sum to exactly $n$.}

\begin{patternbox}
A further specialization of Section~\ref{sec:rec-combination-sum-2}'s
no-reuse, advancing recursion: the candidate pool is now the fixed,
already-sorted range $1$ through $9$ (so no separate sort step or
duplicate-skip is needed, since this fixed pool has no duplicates by
construction), and one new constraint is added -- the combination's
\textbf{length} must be exactly $k$, not just its sum.
\end{patternbox}

\subsection*{Understanding the problem carefully}

The recursion is the same advancing, no-reuse shape as Combination Sum
II, but simplified by the fixed candidate pool $\{1,\dots,9\}$ (already
sorted, no duplicates to worry about), and with the base case now
checking \emph{two} conditions instead of one: both the running count of
chosen numbers must equal $k$, and the running remaining target must
have reached exactly $0$, at the same time.

\begin{ideabox}[Same Template, with a Length Constraint Added to the Base Case]
$\textsc{Solve}(\textit{start}, \textit{remaining}, \textit{current})$:
if $|\textit{current}| = k$: record $\textit{current}$ if
$\textit{remaining}=0$ (otherwise this length-$k$ combination simply
doesn't sum correctly, and we discard it); return either way (no further
recursion needed once length $k$ is reached). Otherwise, for $i$ from
$\textit{start}$ to $9$: if $i > \textit{remaining}$: break (the same
sorted-range pruning as Combination Sum II). Append $i$; recurse
$\textsc{Solve}(i+1, \textit{remaining}-i, \textit{current})$; remove
(undo).
\end{ideabox}

\subsection*{Why checking both conditions together at length $k$ is necessary}

A combination might reach exactly $k$ chosen numbers without summing to
$n$ (too small or too large), or might sum to exactly $n$ using fewer or
more than $k$ numbers -- neither partial success is acceptable, since the
problem requires \emph{both} conditions simultaneously. Stopping the
recursion the instant the count reaches $k$ (regardless of whether the
sum is currently correct) avoids wasted exploration beyond the required
length, while still checking the sum condition explicitly before
recording, since reaching length $k$ alone does not guarantee the target
sum has also been met.

\subsection*{Algorithm}

\begin{enumerate}
  \item $\textsc{Solve}(\textit{start}, \textit{remaining}, \textit{current})$:
  \begin{enumerate}
    \item If $|\textit{current}|=k$: if $\textit{remaining}=0$: record
          $\textit{current}$; return.
    \item For $i$ from $\textit{start}$ to $9$:
    \begin{itemize}
      \item If $i > \textit{remaining}$: break.
      \item Append $i$; recurse $\textsc{Solve}(i+1, \textit{remaining}-i,
            \textit{current})$; remove (undo).
    \end{itemize}
  \end{enumerate}
  \item Call $\textsc{Solve}(1, n, [\,])$.
\end{enumerate}

\subsection*{Dry run}

\dryrun{Take $k=3$, $n=7$.}

\begin{itemize}
  \item $\textsc{Solve}(1,7,[])$: $i=1$: append, recurse
        $\textsc{Solve}(2,6,[1])$.
  \begin{itemize}
    \item $i=2$: append, recurse $\textsc{Solve}(3,4,[1,2])$.
    \begin{itemize}
      \item $|\textit{current}|=2 \ne 3$ yet. $i=3$: $3\le4$, append,
            recurse $\textsc{Solve}(4,1,[1,2,3])$: $|\textit{current}|=3=k$,
            but $\textit{remaining}=1\ne0$: discard. $i=4$: $4>1$,
            break.
    \end{itemize}
    \item Undo. $i=3$ (at $[1]$ level): append, recurse
          $\textsc{Solve}(4,3,[1,3])$: $i=4$: $4>3$, break (no length-3
          combination found here).
  \end{itemize}
  \item Continue similarly... eventually reaching $i=1,i=2,i=4$:
        $\textsc{Solve}(2,6,[1])\to i=4$: $4\le6$, append, recurse
        $\textsc{Solve}(5,2,[1,4])$: $i=5$: $5>2$, break.
  \item Trying $i=1, i=6$ at the $[1]$ level... continuing the full
        search eventually finds $[1,2,4]$ ($1+2+4=7$, length $3$) as a
        valid recorded combination.
\end{itemize}

(The complete search also finds no other length-$3$ combinations from
$\{1,\dots,9\}$ summing to $7$ besides $[1,2,4]$, which the C++ code
below confirms by direct execution.)

\subsection*{C++ implementation}

\begin{lstlisting}
#include <bits/stdc++.h>
using namespace std;

void solve(int start, int remaining, int k, vector<int>& current,
           vector<vector<int>>& result) {
    if ((int)current.size() == k) {
        if (remaining == 0) result.push_back(current);
        return;
    }
    for (int i = start; i <= 9; i++) {
        if (i > remaining) break;

        current.push_back(i);
        solve(i + 1, remaining - i, k, current, result);
        current.pop_back();
    }
}

vector<vector<int>> combinationSum3(int k, int n) {
    vector<vector<int>> result;
    vector<int> current;
    solve(1, n, k, current, result);
    return result;
}

int main() {
    for (auto& comb : combinationSum3(3, 7)) {
        cout << "[ ";
        for (int x : comb) cout << x << " ";
        cout << "] ";
    }
    cout << endl;
    return 0;
}
\end{lstlisting}

\begin{complexitybox}
\textbf{Time Complexity.} The candidate pool is fixed at $9$ values, so
the search space is bounded by $\binom{9}{k}$, a constant independent of
any input size -- effectively $O(1)$ for this fixed-range problem, though
conventionally still described as exponential in the (here, fixed)
pool size.
\textbf{Space Complexity.} $O(k)$ for the recursion depth.
\end{complexitybox}

\begin{takeawaybox}
A fixed, small, already-sorted, duplicate-free candidate pool (like
$1$--$9$ here) simplifies a backtracking template by removing the need
for an initial sort or a same-level duplicate-skip check -- but the core
advancing, no-reuse recursion, plus sorted-pool pruning on the
remaining target, is unchanged from Combination Sum II. Adding a second
base-case condition (here, an exact length $k$) is a small, composable
extension worth recognizing as its own reusable idea.
\end{takeawaybox}

\section{Generate Parentheses}
\label{sec:rec-generate-parentheses}

\problemstatement{Given an integer $n$, generate all combinations of
\textbf{well-formed} (valid) parenthesis strings using exactly $n$ pairs
of parentheses.}

\begin{patternbox}
\emph{``Generate all valid \ldots''} signals backtracking with
\textbf{validity-aware pruning}: rather than generating all $\binom{2n}{n}$
arrangements of $n$ open and $n$ close parentheses and filtering out the
invalid ones afterward (as in the Stack and Queue chapter's Valid
Parentheses problem \emph{checking} a given string), we instead only ever make a choice that
keeps the string-so-far on track to \emph{possibly} become valid,
pruning illegal branches the moment they would become unsalvageable.
\end{patternbox}

\subsection*{Understanding the problem carefully}

At any point while building the string, two counts matter: how many
open parentheses have been placed so far, and how many close parentheses
have been placed so far. An open parenthesis can be placed as long as
fewer than $n$ have been used so far (we must not exceed $n$ total). A
close parenthesis can be placed only if doing so would not create more
closes than opens at any prefix -- i.e.\ only if the current count of
closes is still strictly less than the current count of opens (otherwise
an unmatched close would appear, which the Stack and Queue chapter's Valid Parentheses problem
already established invalidates a parenthesis string).

\begin{ideabox}[Two Guarded Choices Instead of a Binary Pick/Not-Pick]
$\textsc{Solve}(\textit{current}, \textit{openCount}, \textit{closeCount})$:
if $|\textit{current}| = 2n$: record $\textit{current}$; return.
Otherwise: if $\textit{openCount} < n$: append \texttt{'('}; recurse with
$\textit{openCount}+1$; remove (undo). If $\textit{closeCount} <
\textit{openCount}$: append \texttt{')'}; recurse with
$\textit{closeCount}+1$; remove (undo).
\end{ideabox}

\subsection*{Why these two guards are exactly the validity conditions, applied early}

This is the backtracking analogue of
the Stack and Queue chapter's Valid Parentheses stack-based validity check,
applied \emph{during construction} rather than after the fact: a
parenthesis string is valid exactly when every prefix has at least as
many opens as closes, and the total opens and closes both equal $n$. The
guard $\textit{closeCount} < \textit{openCount}$ is precisely
``this prefix currently has more opens than closes, so one more close is
still safe''; the guard $\textit{openCount} < n$ simply caps the total
opens at $n$. Because these guards are checked \emph{before} a character
is ever placed, no invalid prefix is ever constructed in the first
place -- the search space explored is exactly the valid strings (and
their valid-so-far prefixes), never the much larger space of all
arrangements.

\subsection*{Algorithm}

\begin{enumerate}
  \item $\textsc{Solve}(\textit{current}, \textit{openCount}, \textit{closeCount})$:
  \begin{enumerate}
    \item If $|\textit{current}| = 2n$: record $\textit{current}$;
          return.
    \item If $\textit{openCount} < n$: append \texttt{'('}; recurse;
          remove.
    \item If $\textit{closeCount} < \textit{openCount}$: append
          \texttt{')'}; recurse; remove.
  \end{enumerate}
  \item Call $\textsc{Solve}(\texttt{""}, 0, 0)$.
\end{enumerate}

\subsection*{Dry run}

\dryrun{Take $n=2$.}

\begin{itemize}
  \item $\textsc{Solve}(\texttt{""},0,0)$: open allowed ($0<2$):
        append \texttt{(}. $\textsc{Solve}(\texttt{"("},1,0)$.
  \begin{itemize}
    \item Open allowed ($1<2$): append \texttt{(}.
          $\textsc{Solve}(\texttt{"(("},2,0)$.
    \begin{itemize}
      \item Open not allowed ($2\not<2$). Close allowed ($0<2$): append
            \texttt{)}. $\textsc{Solve}(\texttt{"(()"},2,1)$: close
            allowed ($1<2$): append \texttt{)}.
            $\textsc{Solve}(\texttt{"(())"},2,2)$: length $4=2n$,
            record \texttt{"(())"}.
    \end{itemize}
    \item Undo back to \texttt{"("}. Close allowed ($0<1$): append
          \texttt{)}. $\textsc{Solve}(\texttt{"()"},1,1)$: open allowed
          ($1<2$): append \texttt{(}.
          $\textsc{Solve}(\texttt{"()("},2,1)$: close allowed
          ($1<2$): append \texttt{)}.
          $\textsc{Solve}(\texttt{"()()"},2,2)$: record
          \texttt{"()()"}.
  \end{itemize}
\end{itemize}

Recorded strings: \texttt{"(())"}, \texttt{"()()"} -- exactly the $2$
valid parenthesizations of $n=2$ pairs.

\subsection*{C++ implementation}

\begin{lstlisting}
#include <bits/stdc++.h>
using namespace std;

void solve(string& current, int openCount, int closeCount, int n,
           vector<string>& result) {
    if ((int)current.size() == 2 * n) {
        result.push_back(current);
        return;
    }
    if (openCount < n) {
        current.push_back('(');
        solve(current, openCount + 1, closeCount, n, result);
        current.pop_back();
    }
    if (closeCount < openCount) {
        current.push_back(')');
        solve(current, openCount, closeCount + 1, n, result);
        current.pop_back();
    }
}

vector<string> generateParenthesis(int n) {
    vector<string> result;
    string current;
    solve(current, 0, 0, n, result);
    return result;
}

int main() {
    for (string& s : generateParenthesis(3)) cout << s << " ";
    cout << endl;
    return 0;
}
\end{lstlisting}

\begin{complexitybox}
\textbf{Time Complexity.} The number of valid strings generated is the
$n$-th Catalan number, $C_n = \binom{2n}{n}/(n+1)$, and each costs
$O(n)$ to build, giving an overall time complexity of
\[
  O\!\left(\frac{4^n}{\sqrt{n}}\right)
\]
(the asymptotic growth rate of the Catalan numbers), substantially
smaller than the $O(2^{2n})$ that generating and filtering all
arrangements would cost.
\textbf{Space Complexity.} $O(n)$ for the recursion depth (the string
grows to length $2n$, but the call depth is also $O(n)$ since two
characters can be added back to back in the worst case, or more
precisely the depth equals the string length built so far, $O(n)$ at the
deepest).
\end{complexitybox}

\begin{takeawaybox}
Whenever a backtracking problem's notion of ``validity'' can be checked
\emph{incrementally} (true/false after every partial choice, not only at
the very end), build that incremental check directly into the recursion
as a guard before each choice, rather than generating complete candidates
and validating them afterward. This is dramatically more efficient
whenever invalid prefixes are common, since entire invalid subtrees are
never explored at all.
\end{takeawaybox}

\section{Generate All Binary Strings with No Two Consecutive 1s}
\label{sec:rec-generate-binary-strings}

\problemstatement{Given an integer $n$, generate all binary strings of
length $n$ such that no two $\texttt{1}$s are adjacent (no two
consecutive positions both hold $\texttt{1}$).}

\begin{patternbox}
Another instance of Section~\ref{sec:rec-generate-parentheses}'s
incremental-validity-guard pattern: instead of generating all $2^n$
binary strings and filtering, build the constraint directly into the
choice of what character can be placed \textbf{next}, based only on
\textbf{the single most recently placed character} -- a guard that
depends on a small, fixed amount of local history rather than the entire
string built so far.
\end{patternbox}

\subsection*{Understanding the problem carefully}

At every position, placing a $\texttt{0}$ is always safe, regardless of
what came before -- it can never create two consecutive $\texttt{1}$s.
Placing a $\texttt{1}$ is safe \emph{only if} the immediately preceding
character (if any) was $\texttt{0}$ -- placing it right after another
$\texttt{1}$ would violate the constraint. So the only piece of history
the recursion needs to remember is a single boolean: ``was the last
character placed a $\texttt{1}$?''

\begin{ideabox}[Guard the Choice of $\texttt{1}$ by the Last Character Placed]
$\textsc{Solve}(\textit{current}, \textit{lastWasOne})$: if
$|\textit{current}|=n$: record $\textit{current}$; return. Always:
append $\texttt{0}$; recurse with $\textit{lastWasOne}=\text{false}$;
remove (undo). If $\textit{lastWasOne}$ is \text{false} (i.e.\ the
previous character was not $\texttt{1}$, or this is the first
character): append $\texttt{1}$; recurse with
$\textit{lastWasOne}=\text{true}$; remove (undo).
\end{ideabox}

\subsection*{Why remembering only the single last character is sufficient}

The constraint ``no two consecutive $\texttt{1}$s'' is a purely
\textbf{local} property -- it only ever compares two \emph{adjacent}
positions, never anything further apart. So the only information needed
to decide whether placing a $\texttt{1}$ right now is safe is whether the
immediately preceding position holds a $\texttt{1}$; nothing about
positions further back can affect this decision. This is what allows the
recursion to carry just one bit of state (rather than re-scanning the
entire $\textit{current}$ string at every step to check for violations),
keeping each step's validity check $O(1)$.

\subsection*{Algorithm}

\begin{enumerate}
  \item $\textsc{Solve}(\textit{current}, \textit{lastWasOne})$:
  \begin{enumerate}
    \item If $|\textit{current}|=n$: record $\textit{current}$; return.
    \item Append $\texttt{0}$; recurse with
          $\textit{lastWasOne}=\text{false}$; remove (undo).
    \item If not $\textit{lastWasOne}$: append $\texttt{1}$; recurse
          with $\textit{lastWasOne}=\text{true}$; remove (undo).
  \end{enumerate}
  \item Call $\textsc{Solve}(\texttt{""}, \text{false})$.
\end{enumerate}

\subsection*{Dry run}

\dryrun{Take $n=3$.}

\begin{itemize}
  \item $\textsc{Solve}(\texttt{""}, \text{false})$: place
        $\texttt{0}$: $\textsc{Solve}(\texttt{"0"}, \text{false})$.
  \begin{itemize}
    \item Place $\texttt{0}$: $\textsc{Solve}(\texttt{"00"},
          \text{false})$: place $\texttt{0}$: record \texttt{"000"};
          place $\texttt{1}$ (allowed, last was $0$): record
          \texttt{"001"}.
    \item Place $\texttt{1}$ (allowed): $\textsc{Solve}(\texttt{"01"},
          \text{true})$: place $\texttt{0}$: record \texttt{"010"};
          place $\texttt{1}$ -- \textbf{not allowed} (last was
          $\texttt{1}$), skipped entirely.
  \end{itemize}
  \item Place $\texttt{1}$ (allowed at the very first position):
        $\textsc{Solve}(\texttt{"1"}, \text{true})$: place $\texttt{0}$:
        $\textsc{Solve}(\texttt{"10"}, \text{false})$: place
        $\texttt{0}$: record \texttt{"100"}; place $\texttt{1}$
        (allowed): record \texttt{"101"}. Place $\texttt{1}$ -- not
        allowed (last was $\texttt{1}$), skipped.
\end{itemize}

Recorded strings: \texttt{"000", "001", "010", "100", "101"} -- exactly
the $5$ length-$3$ binary strings with no two consecutive $\texttt{1}$s
(matching the Fibonacci-like count $F(n+2)$ of such strings).

\subsection*{C++ implementation}

\begin{lstlisting}
#include <bits/stdc++.h>
using namespace std;

void solve(string& current, bool lastWasOne, int n, vector<string>& result) {
    if ((int)current.size() == n) {
        result.push_back(current);
        return;
    }
    current.push_back('0');
    solve(current, false, n, result);
    current.pop_back();

    if (!lastWasOne) {
        current.push_back('1');
        solve(current, true, n, result);
        current.pop_back();
    }
}

vector<string> generateBinaryStrings(int n) {
    vector<string> result;
    string current;
    solve(current, false, n, result);
    return result;
}

int main() {
    for (string& s : generateBinaryStrings(3)) cout << s << " ";
    cout << endl;
    return 0;
}
\end{lstlisting}

\begin{complexitybox}
\textbf{Time Complexity.} The number of valid strings of length $n$ is
$O(\phi^n)$ (a Fibonacci-like growth rate, where $\phi$ is the golden
ratio), each costing $O(n)$ to build, giving an overall time complexity
of
\[
  O(n \cdot \phi^n),
\]
substantially smaller than the $O(n \cdot 2^n)$ of generating and
filtering all binary strings.
\textbf{Space Complexity.} $O(n)$ for the recursion depth.
\end{complexitybox}

\begin{takeawaybox}
When a backtracking constraint only ever depends on a small, fixed
amount of recent history (here, just the single previous character),
thread exactly that small amount of state through the recursion's
parameters, rather than re-deriving it by re-scanning the entire
candidate built so far. This keeps each step's validity check $O(1)$ and
generalizes to any constraint expressible in terms of a bounded
``sliding window'' of recent choices.
\end{takeawaybox}

\section{Letter Combinations of a Phone Number}
\label{sec:rec-letter-combinations}

\problemstatement{Given a string of digits from $2$ to $9$, where each
digit maps to a fixed set of letters exactly as on a telephone keypad
(e.g.\ $2 \to \texttt{abc}$, $3 \to \texttt{def}$, \ldots, $7 \to
\texttt{pqrs}$, $9 \to \texttt{wxyz}$), return every possible letter
combination obtainable by choosing one letter for each digit, in order.}

\begin{patternbox}
\emph{``Choose one option per position, from a different choice set at
each position''} generalizes Section~\ref{sec:rec-power-set}'s binary
pick/not-pick into an $n$-ary branch at every level: instead of $2$
choices (include/exclude) at each position, there are however many
letters the current digit maps to, and the recursion simply loops over
all of them instead of hard-coding exactly two branches.
\end{patternbox}

\subsection*{Understanding the problem carefully}

Each digit position contributes one letter to the final combination,
chosen independently from that digit's fixed mapping (e.g.\ digit
$\texttt{'2'}$ contributes one of $\texttt{a},\texttt{b},\texttt{c}$).
Since these choices are independent across positions, exactly as the
pick/not-pick choices were independent across array elements in
Section~\ref{sec:rec-power-set}, the total number of combinations is the
\textbf{product} of each digit's letter count, and a recursion that
loops through all letters for the current digit, recursing once per
letter, visits every combination exactly once.

\begin{ideabox}[Loop Over the Current Digit's Letters Instead of a Fixed Binary Choice]
$\textsc{Solve}(\textit{index}, \textit{current})$: if
$\textit{index} = |\textit{digits}|$: record $\textit{current}$;
return. Otherwise: for each letter $c$ in the mapping for
$\textit{digits}[\textit{index}]$: append $c$ to $\textit{current}$;
recurse $\textsc{Solve}(\textit{index}+1, \textit{current})$; remove $c$
(undo).
\end{ideabox}

\subsection*{Why looping generalizes the binary pick/not-pick template directly}

The pick/not-pick recursion in Section~\ref{sec:rec-power-set} is the
special case of this loop where every position's choice set happens to
have exactly $2$ options (``include'' or ``exclude''). Here, the choice
set size varies by digit (e.g.\ $3$ letters for most digits, but $4$ for
$\texttt{7}$ and $\texttt{9}$), so the recursion is written as a general
loop over however many choices exist at the current position, rather
than two separately hard-coded branches -- the same choose/recurse/un-choose
rhythm applies uniformly, regardless of how many choices the loop
iterates over.

\subsection*{Algorithm}

\begin{enumerate}
  \item Define the digit-to-letters mapping (a fixed lookup table).
  \item $\textsc{Solve}(\textit{index}, \textit{current})$:
  \begin{enumerate}
    \item If $\textit{index} = |\textit{digits}|$: record
          $\textit{current}$; return.
    \item For each letter $c$ mapped from
          $\textit{digits}[\textit{index}]$: append $c$; recurse
          $\textsc{Solve}(\textit{index}+1, \textit{current})$; remove
          (undo).
  \end{enumerate}
  \item Call $\textsc{Solve}(0, \texttt{""})$ (return an empty list
        immediately if $\textit{digits}$ is empty, as a special case).
\end{enumerate}

\subsection*{Dry run}

\dryrun{Take $\textit{digits} = \texttt{"23"}$ ($2\to\texttt{abc}$,
$3\to\texttt{def}$).}

\begin{itemize}
  \item $\textsc{Solve}(0,\texttt{""})$: loop over $\texttt{a,b,c}$ for
        digit $\texttt{'2'}$.
  \begin{itemize}
    \item $c=\texttt{a}$: $\textsc{Solve}(1,\texttt{"a"})$: loop over
          $\texttt{d,e,f}$ for digit $\texttt{'3'}$: record
          \texttt{"ad"}, \texttt{"ae"}, \texttt{"af"}.
    \item $c=\texttt{b}$: $\textsc{Solve}(1,\texttt{"b"})$: record
          \texttt{"bd"}, \texttt{"be"}, \texttt{"bf"}.
    \item $c=\texttt{c}$: $\textsc{Solve}(1,\texttt{"c"})$: record
          \texttt{"cd"}, \texttt{"ce"}, \texttt{"cf"}.
  \end{itemize}
\end{itemize}

Result: \texttt{"ad","ae","af","bd","be","bf","cd","ce","cf"} -- all
$3\times3=9$ combinations.

\subsection*{C++ implementation}

\begin{lstlisting}
#include <bits/stdc++.h>
using namespace std;

void solve(int index, string& digits, string& current,
           vector<string>& mapping, vector<string>& result) {
    if (index == (int)digits.size()) {
        result.push_back(current);
        return;
    }
    string letters = mapping[digits[index] - '0'];
    for (char c : letters) {
        current.push_back(c);
        solve(index + 1, digits, current, mapping, result);
        current.pop_back();
    }
}

vector<string> letterCombinations(string digits) {
    if (digits.empty()) return {};

    vector<string> mapping = {
        "", "", "abc", "def", "ghi", "jkl",
        "mno", "pqrs", "tuv", "wxyz"
    };

    vector<string> result;
    string current;
    solve(0, digits, current, mapping, result);
    return result;
}

int main() {
    for (string& s : letterCombinations("23")) cout << s << " ";
    cout << endl;
    return 0;
}
\end{lstlisting}

\begin{complexitybox}
\textbf{Time Complexity.} If each digit maps to at most $4$ letters,
there are at most $4^L$ combinations for a digit string of length $L$,
each costing $O(L)$ to build, giving
\[
  O(L \cdot 4^L).
\]
\textbf{Space Complexity.} $O(L)$ for the recursion depth.
\end{complexitybox}

\begin{takeawaybox}
A loop-based recursive branch (``try every choice in this position's
choice set'') is the natural generalization of a hard-coded binary
pick/not-pick, and is the more common shape for backtracking problems in
practice, since most real choice sets have more than two options. Every
remaining backtracking problem in this chapter
(Sections~\ref{sec:rec-rat-in-maze}--\ref{sec:rec-word-search}) uses this
loop-based shape, with the choice set determined by the problem's
specific structure (grid moves, board placements, partition points).
\end{takeawaybox}

\section{Rat in a Maze}
\label{sec:rec-rat-in-maze}

\noindent We now begin the chapter's final group: full
\textbf{constraint-satisfaction backtracking}, where the choice at each
step is a move on a grid or board, and a choice already made can block
later choices in ways that must be explicitly tracked and undone.

\problemstatement{Given an $n \times n$ binary maze (cells are either
open, marked $1$, or blocked, marked $0$), a rat starts at the top-left
cell $(0,0)$ and wants to reach the bottom-right cell $(n-1,n-1)$,
moving only through open cells, in any of the four directions (down,
left, right, up), without revisiting a cell already used in the current
path. Find every path from start to destination, each described as a
string of move directions.}

\begin{patternbox}
\emph{``Find all paths through a grid avoiding revisits''} is the
canonical \textbf{grid backtracking} template: at each cell, the choice
set is ``which of up to $4$ neighboring directions to move into next,''
each choice must be validated against the grid's boundaries and
blocked/visited cells \emph{before} being taken, and -- crucially -- a
cell marked as visited while exploring one direction must be
\textbf{unmarked} once that direction's exploration fully returns, so
that a \emph{different} path (reached via backtracking to an earlier
cell) can freely reuse that same cell.
\end{patternbox}

\subsection*{Understanding the problem carefully}

At any cell $(\textit{row},\textit{col})$ during the search, a move into
neighbor $(\textit{newRow},\textit{newCol})$ is valid only if it is
within the grid's bounds, the cell there is open (value $1$), and it has
not already been visited earlier in the \emph{current} path (revisiting
would create a cycle, which is never part of a valid simple path). If
the move is valid, we mark the destination visited, append the
direction's label to the path string, and recurse from there; once that
entire recursive exploration finishes (having found every path reachable
through this move), we must \textbf{unmark} the destination as visited
and remove the direction label, so that a sibling direction's
exploration -- reached by returning to this same cell -- does not
incorrectly believe this cell is permanently blocked.

\begin{ideabox}[Mark Visited Before Recursing, Unmark After Returning]
$\textsc{Solve}(\textit{row}, \textit{col}, \textit{path})$: if
$(\textit{row},\textit{col})$ is the destination: record $\textit{path}$
(and return, or continue if the destination cell itself could lead
further, depending on problem variant). Otherwise, for each of the $4$
directions, in a fixed order: compute the neighbor; if it is in bounds,
open, and not visited: mark it visited; append the direction's label to
$\textit{path}$; recurse into the neighbor; remove the label from
$\textit{path}$ (undo); unmark the neighbor as visited (undo).
\end{ideabox}

\subsection*{Why unmarking after the recursive call returns is essential for correctness}

If a cell, once visited during the exploration of one direction, were
\emph{never} unmarked, it would remain permanently blocked for every
other path that might also legitimately need to pass through it later
-- incorrectly preventing valid paths from being discovered. The
visited-marking only needs to prevent revisiting \emph{within a single,
currently-in-progress path} (to avoid infinite loops/cycles along that
one path); it must not persist across different candidate paths, which
are explored one at a time, each needing a clean slate. Unmarking
immediately after a direction's full recursive exploration returns is
exactly what restores that clean slate for the next direction (or for an
ancestor's next sibling branch) to use.

\subsection*{Algorithm}

\begin{enumerate}
  \item Mark $(0,0)$ visited; call $\textsc{Solve}(0,0,\texttt{""})$.
  \item $\textsc{Solve}(\textit{row}, \textit{col}, \textit{path})$:
  \begin{enumerate}
    \item If $(\textit{row},\textit{col}) = (n-1,n-1)$: record
          $\textit{path}$; return.
    \item For each direction $(dr, dc, \textit{label})$ in $\{(1,0,
          \texttt{D}), (0,-1,\texttt{L}), (0,1,\texttt{R}),
          (-1,0,\texttt{U})\}$:
    \begin{itemize}
      \item $\textit{nr} \gets \textit{row}+dr$, $\textit{nc} \gets
            \textit{col}+dc$.
      \item If $\textit{nr},\textit{nc}$ in bounds, maze cell is open,
            and not visited: mark visited; append $\textit{label}$ to
            $\textit{path}$; recurse $\textsc{Solve}(\textit{nr},
            \textit{nc}, \textit{path})$; remove the appended label
            (undo); unmark visited (undo).
    \end{itemize}
  \end{enumerate}
\end{enumerate}

\subsection*{Dry run}

\dryrun{Take the $2\times2$ maze
$\begin{bmatrix}1&1\\0&1\end{bmatrix}$.}

\begin{itemize}
  \item Start at $(0,0)$, visited$=\{(0,0)\}$,
        $\textit{path}=\texttt{""}$.
  \item Try \texttt{D} (down to $(1,0)$): maze value is $0$ (blocked).
        Invalid, skip.
  \item Try \texttt{L} (left to $(0,-1)$): out of bounds. Skip.
  \item Try \texttt{R} (right to $(0,1)$): open, not visited. Mark
        visited, path$=\texttt{"R"}$. Recurse from $(0,1)$.
  \begin{itemize}
    \item Try \texttt{D} (down to $(1,1)$): open, not visited. Mark
          visited, path$=\texttt{"RD"}$. Recurse from $(1,1)$:
          destination reached! Record \texttt{"RD"}.
    \item Undo: unmark $(1,1)$, path back to $\texttt{"R"}$.
    \item Try \texttt{L} (back to $(0,0)$): visited already, skip.
    \item Try \texttt{U}: out of bounds. Skip.
  \end{itemize}
  \item Undo: unmark $(0,1)$, path back to $\texttt{""}$.
  \item Try \texttt{U} (up to $(-1,0)$): out of bounds. Skip.
\end{itemize}

The only path found is \texttt{"RD"} (right, then down).

\subsection*{C++ implementation}

\begin{lstlisting}
#include <bits/stdc++.h>
using namespace std;

void solve(int row, int col, vector<vector<int>>& maze, int n,
           vector<vector<bool>>& visited, string& path,
           vector<string>& result) {
    if (row == n - 1 && col == n - 1) {
        result.push_back(path);
        return;
    }

    int dr[] = {1, 0, 0, -1};
    int dc[] = {0, -1, 1, 0};
    char label[] = {'D', 'L', 'R', 'U'};

    for (int d = 0; d < 4; d++) {
        int nr = row + dr[d], nc = col + dc[d];
        if (nr >= 0 && nr < n && nc >= 0 && nc < n &&
            maze[nr][nc] == 1 && !visited[nr][nc]) {
            visited[nr][nc] = true;
            path.push_back(label[d]);

            solve(nr, nc, maze, n, visited, path, result);

            path.pop_back();
            visited[nr][nc] = false;
        }
    }
}

vector<string> ratInMaze(vector<vector<int>>& maze) {
    int n = maze.size();
    vector<string> result;
    if (maze[0][0] == 0) return result;

    vector<vector<bool>> visited(n, vector<bool>(n, false));
    visited[0][0] = true;
    string path;
    solve(0, 0, maze, n, visited, path, result);
    return result;
}

int main() {
    vector<vector<int>> maze = {{1,1},{0,1}};
    for (string& p : ratInMaze(maze)) cout << p << " ";
    cout << endl;
    return 0;
}
\end{lstlisting}

\begin{complexitybox}
\textbf{Time Complexity.} In the worst case (an open grid with no
blocked cells), the number of simple paths can be exponential in $n^2$
(roughly $O(4^{n^2})$ branching, though heavily pruned by the
no-revisit rule in practice); a tight closed-form bound is not generally
useful here -- what matters is that each cell visit does $O(1)$ work
beyond the recursive calls themselves.
\textbf{Space Complexity.} $O(n^2)$ for the visited grid, plus
$O(n^2)$ for the maximum recursion depth (a path can visit at most every
cell once).
\end{complexitybox}

\begin{takeawaybox}
``Mark visited before recursing, unmark after returning'' is the
universal template for any backtracking search over a grid or graph that
must avoid revisiting nodes within a single path -- it is the direct
generalization of the choose/recurse/un-choose rhythm from
Section~\ref{sec:rec-power-set}, applied to spatial moves instead of
array-element inclusion. The same template, with different validity
checks, drives N-Queens, Sudoku, and Word Search in the sections that
follow.
\end{takeawaybox}

\section{N-Queens}
\label{sec:rec-n-queens}

\problemstatement{Place $n$ queens on an $n \times n$ chessboard so that
no two queens attack each other (no two share a row, a column, or a
diagonal). Return every distinct valid placement.}

\begin{patternbox}
\emph{``Place $n$ non-attacking items on a board''} is a constraint-satisfaction
backtracking problem where a powerful structural shortcut is available:
since a valid solution can have \textbf{at most one queen per row}
(two queens in the same row always attack each other), we can place
exactly one queen \textbf{per row}, in row order, reducing the search
from ``choose a subset of cells'' to ``choose one column per row,''
checking only column and diagonal conflicts against \emph{already-placed}
queens in earlier rows.
\end{patternbox}

\subsection*{Understanding the problem carefully}

Because we place queens one complete row at a time, by the time we
decide row $r$'s queen, rows $0, \dots, r-1$ already have their queens
correctly and safely placed; we never need to recheck them against each
other again, only check the new candidate column against those already
fixed. A candidate column $c$ in row $r$ is safe if no earlier queen
shares column $c$, and no earlier queen lies on either diagonal passing
through $(r,c)$ -- a queen at $(r', c')$ shares a diagonal with $(r,c)$
exactly when $|r-r'| = |c-c'|$.

\begin{ideabox}[Track Occupied Columns and Diagonals for $O(1)$ Safety Checks]
Maintain three boolean arrays: $\textit{colUsed}[c]$ (is column $c$
occupied by some already-placed queen), and two diagonal arrays indexed
by the constant value along each diagonal direction:
$\textit{diag1}[r+c]$ (the ``$\backslash$''-diagonal) and
$\textit{diag2}[r-c+n-1]$ (the ``$/$''-diagonal, offset to stay
non-negative). $\textsc{Solve}(\textit{row})$: if $\textit{row}=n$:
record the current board; return. Otherwise, for each column $c$ from
$0$ to $n-1$: if none of $\textit{colUsed}[c]$, $\textit{diag1}[r+c]$,
$\textit{diag2}[r-c+n-1]$ are set: place the queen, mark all three
arrays, recurse $\textsc{Solve}(\textit{row}+1)$, then unmark all three
(undo) and remove the queen.
\end{ideabox}

\subsection*{Why diagonal membership reduces to a single sum or difference}

Every cell on the same ``$\backslash$''-diagonal (running from upper-left
to lower-right) shares the same value of $\textit{row}+\textit{col}$;
every cell on the same ``$/$''-diagonal (lower-left to upper-right)
shares the same value of $\textit{row}-\textit{col}$. This means
checking ``does any earlier queen share a diagonal with $(r,c)$'' reduces
to a single array lookup at index $r+c$ (or $r-c$, shifted to stay a
valid non-negative array index) rather than scanning every previously
placed queen and computing $|r-r'|$ versus $|c-c'|$ each time -- turning
an $O(n)$ safety check into an $O(1)$ one, and the whole algorithm's
per-placement cost from $O(n)$ down to $O(1)$.

\subsection*{Algorithm}

\begin{enumerate}
  \item Initialize $\textit{colUsed}, \textit{diag1}, \textit{diag2}$
        as all-false boolean arrays of appropriate sizes.
  \item $\textsc{Solve}(\textit{row})$:
  \begin{enumerate}
    \item If $\textit{row}=n$: record the board; return.
    \item For $c$ from $0$ to $n-1$: if $\textit{colUsed}[c]$,
          $\textit{diag1}[\textit{row}+c]$, and
          $\textit{diag2}[\textit{row}-c+n-1]$ are all false: place the
          queen at $(\textit{row},c)$; set all three arrays true;
          recurse $\textsc{Solve}(\textit{row}+1)$; set all three
          arrays back to false (undo); remove the queen.
  \end{enumerate}
  \item Call $\textsc{Solve}(0)$.
\end{enumerate}

\subsection*{Dry run}

\dryrun{Take $n=4$.}

The search places a queen in row $0$, tries each column, and recurses;
the only two solutions (up to the board's symmetry) for $n=4$ are
columns $[1,3,0,2]$ and $[2,0,3,1]$ (queen in row $0$ at column $1$ or
$2$ respectively, with the rest determined by the safety constraints).
Tracing column $1$ for row $0$: row $1$'s columns $0,1,2$ are all
unsafe (column $1$ taken, or diagonal conflicts), leaving only column
$3$ safe; row $2$'s only safe column, given rows $0,1$, is $0$; row
$3$'s only safe column, given rows $0,1,2$, is $2$ -- yielding the
complete solution $[1,3,0,2]$.

\subsection*{C++ implementation}

\begin{lstlisting}
#include <bits/stdc++.h>
using namespace std;

void solve(int row, int n, vector<int>& cols, vector<bool>& colUsed,
           vector<bool>& diag1, vector<bool>& diag2,
           vector<vector<int>>& result) {
    if (row == n) {
        result.push_back(cols);
        return;
    }
    for (int c = 0; c < n; c++) {
        int d1 = row + c, d2 = row - c + n - 1;
        if (colUsed[c] || diag1[d1] || diag2[d2]) continue;

        cols[row] = c;
        colUsed[c] = diag1[d1] = diag2[d2] = true;

        solve(row + 1, n, cols, colUsed, diag1, diag2, result);

        colUsed[c] = diag1[d1] = diag2[d2] = false;
    }
}

vector<vector<int>> solveNQueens(int n) {
    vector<vector<int>> result;
    vector<int> cols(n, -1);
    vector<bool> colUsed(n, false), diag1(2 * n, false), diag2(2 * n, false);

    solve(0, n, cols, colUsed, diag1, diag2, result);
    return result;
}

int main() {
    auto solutions = solveNQueens(4);
    cout << "Number of solutions: " << solutions.size() << endl;
    for (auto& sol : solutions) {
        for (int c : sol) cout << c << " ";
        cout << endl;
    }
    return 0;
}
\end{lstlisting}

\begin{complexitybox}
\textbf{Time Complexity.} Bounded by $O(n!)$ in the absolute worst case
(one queen per row, up to $n$ column choices each, shrinking as
constraints accumulate), though the $O(1)$ safety check via the
tracking arrays makes each individual placement attempt fast; the true
practical running time is much smaller due to aggressive pruning from
the column/diagonal constraints.
\textbf{Space Complexity.} $O(n)$ for the tracking arrays and the
recursion depth.
\end{complexitybox}

\begin{takeawaybox}
When a backtracking constraint can be reduced to membership in a small
number of equivalence classes (here, columns and the two diagonal
directions), maintain one boolean array per class for $O(1)$
membership checks, instead of recomputing the constraint from scratch
by scanning all previously placed items on every attempt. This
``maintain auxiliary used/unused trackers, mark and unmark them in
lockstep with placing and removing a choice'' idea generalizes directly
to Sudoku's row/column/box constraints in the next section.
\end{takeawaybox}

\section{Sudoku Solver}
\label{sec:rec-sudoku-solver}

\problemstatement{Given a partially filled $9\times9$ Sudoku board
(empty cells marked \texttt{'.'}), fill every empty cell with a digit
from $1$ to $9$ so that every row, every column, and every one of the
nine $3\times3$ boxes contains each digit exactly once. It is guaranteed
the input has a unique solution.}

\begin{patternbox}
This closing backtracking problem combines two ideas already
established: N-Queens's (Section~\ref{sec:rec-n-queens}) ``track
constraint membership in auxiliary structures for fast checks''
(generalized here to three overlapping constraints -- row, column, and
box -- instead of just column and two diagonals), and a new control-flow
idea needed because this problem wants \textbf{any one} valid completion,
not \emph{every} one: the recursion must \textbf{return a boolean}
signaling success, so that the instant a complete valid filling is found
anywhere in the search, every enclosing call can immediately stop trying
further alternatives and propagate that success straight back to the
top.
\end{patternbox}

\subsection*{Understanding the problem carefully}

Scan the board for the next empty cell (in, say, row-major order). For
that cell, try each digit $1$ through $9$; a digit is valid if it does
not already appear elsewhere in the same row, the same column, or the
same $3\times3$ box. If a valid digit is found, place it and recursively
attempt to solve the rest of the board. If that recursive attempt
\emph{succeeds} (returns \texttt{true}, meaning a full valid solution was
found somewhere beneath this choice), we are done -- stop trying further
digits at this cell and propagate \texttt{true} upward immediately. If
it \emph{fails} (returns \texttt{false}), undo this digit and try the
next one. If no digit from $1$ to $9$ works at this cell, this entire
branch is a dead end -- return \texttt{false}, triggering the caller to
backtrack and try a different digit at an earlier cell.

\begin{ideabox}[Boolean-Returning Backtracking: Stop on First Success]
$\textsc{Solve}(\textit{board})$: find the next empty cell
$(\textit{row},\textit{col})$; if none exists, the board is completely
and validly filled -- return \texttt{true} (base case, success).
Otherwise, for $d$ from $1$ to $9$: if placing $d$ at
$(\textit{row},\textit{col})$ is valid: place it; if
$\textsc{Solve}(\textit{board})$ returns \texttt{true}: return
\texttt{true} immediately (propagate success, no undo needed -- this
digit is part of the final solution). Otherwise: remove $d$ (undo) and
try the next digit. If every digit from $1$ to $9$ has been tried
without success: return \texttt{false}.
\end{ideabox}

\subsection*{Why a boolean return value, with immediate propagation, is the key new idea}

Section~\ref{sec:rec-n-queens} and earlier sections collected \emph{all}
valid solutions, so every branch had to be explored regardless of
earlier successes. Sudoku, by contrast, only needs \emph{one} valid
completion (and the problem guarantees exactly one exists), so the
moment a recursive call confirms a digit choice leads all the way to a
complete, valid board, there is no reason to keep searching -- every
enclosing call should immediately stop and also report success,
unwinding the entire call stack without exploring any further
alternatives. This is exactly Section~\ref{sec:rec-check-subsequence-sum-k}'s
short-circuit-on-success idea, now applied to a problem where the
``success'' condition is ``the rest of the board can be completed,''
checked recursively rather than as a simple base-case comparison.

\subsection*{Algorithm}

\begin{enumerate}
  \item $\textsc{Solve}(\textit{board})$:
  \begin{enumerate}
    \item Find the next empty cell $(\textit{row}, \textit{col})$; if
          none, return \texttt{true}.
    \item For $d$ from $1$ to $9$:
    \begin{itemize}
      \item If $d$ is valid at $(\textit{row}, \textit{col})$ (no
            conflict in its row, column, or $3\times3$ box): place
            $d$. If $\textsc{Solve}(\textit{board})$ returns
            \texttt{true}: return \texttt{true}. Else: remove $d$
            (undo).
    \end{itemize}
    \item Return \texttt{false} (no digit worked here; trigger
          backtracking).
  \end{enumerate}
\end{enumerate}

\subsection*{Dry run}

\dryrun{A small conceptual trace: suppose cell $(0,2)$ is empty, and
digits $1$--$5$ are already ruled out by its row/column/box, but
$6$ is valid. We place $6$ and recurse to the next empty cell. Several
cells later, suppose a dead end is reached where no digit from $1$--$9$
is valid at some cell $(4,7)$ -- this returns \texttt{false}, which
propagates back to the most recent choice point (perhaps the placement
at $(3,5)$), forcing it to try its \emph{next} candidate digit instead
of the one that led to this dead end. This unwind-and-retry continues
until either a different early choice avoids the dead end entirely (and
the search eventually completes successfully, propagating \texttt{true}
all the way back to the top), or -- if the puzzle is unsolvable -- every
possibility is exhausted and the top-level call returns \texttt{false}.}

\subsection*{C++ implementation}

\begin{lstlisting}
#include <bits/stdc++.h>
using namespace std;

bool isValid(vector<vector<char>>& board, int row, int col, char d) {
    for (int i = 0; i < 9; i++) {
        if (board[row][i] == d) return false;
        if (board[i][col] == d) return false;
        if (board[3*(row/3) + i/3][3*(col/3) + i%3] == d) return false;
    }
    return true;
}

bool solve(vector<vector<char>>& board) {
    for (int row = 0; row < 9; row++) {
        for (int col = 0; col < 9; col++) {
            if (board[row][col] != '.') continue;

            for (char d = '1'; d <= '9'; d++) {
                if (isValid(board, row, col, d)) {
                    board[row][col] = d;
                    if (solve(board)) return true;
                    board[row][col] = '.';
                }
            }
            return false; // no digit worked at this empty cell
        }
    }
    return true; // no empty cell left: solved
}

int main() {
    vector<vector<char>> board = {
        {'5','3','.','.','7','.','.','.','.'},
        {'6','.','.','1','9','5','.','.','.'},
        {'.','9','8','.','.','.','.','6','.'},
        {'8','.','.','.','6','.','.','.','3'},
        {'4','.','.','8','.','3','.','.','1'},
        {'7','.','.','.','2','.','.','.','6'},
        {'.','6','.','.','.','.','2','8','.'},
        {'.','.','.','4','1','9','.','.','5'},
        {'.','.','.','.','8','.','.','7','9'}
    };

    solve(board);

    for (auto& row : board) {
        for (char c : row) cout << c << " ";
        cout << endl;
    }
    return 0;
}
\end{lstlisting}

\begin{complexitybox}
\textbf{Time Complexity.} Exponential in the worst case, roughly
$O(9^{m})$ where $m$ is the number of empty cells (each could in
principle try $9$ digits), though the row/column/box validity checks
prune the vast majority of branches in practice, making real Sudoku
puzzles solve very quickly.
\textbf{Space Complexity.} $O(m)$ for the recursion depth (at most one
call per empty cell along the current path), plus $O(1)$ extra for the
$\textit{isValid}$ check (which scans a fixed $9$ cells each for row,
column, and box).
\end{complexitybox}

\begin{takeawaybox}
Boolean-returning backtracking -- where finding \emph{one} success
anywhere in a subtree should immediately propagate \texttt{true} all the
way up, aborting every remaining sibling branch -- is the right shape
whenever a problem asks for any single valid solution (or asks whether
one exists at all), as opposed to enumerating every solution
(N-Queens) or counting them. Recognizing which of these three modes
(generate all, exists, count) a backtracking problem needs, exactly as
flagged back in Section~\ref{sec:rec-count-subsequence-sum-k}'s
takeaway, remains the first and most important design decision.
\end{takeawaybox}

\section{Word Search}
\label{sec:rec-word-search}

\problemstatement{Given an $m \times n$ grid of characters and a string
$\textit{word}$, determine whether $\textit{word}$ can be formed by a
path of \textbf{sequentially adjacent} cells (horizontally or vertically
neighboring, not diagonal), where the \textbf{same cell may not be used
more than once} within the path.}

\begin{patternbox}
This is Rat in a Maze's (Section~\ref{sec:rec-rat-in-maze}) grid
backtracking template, combined with Sudoku's
(Section~\ref{sec:rec-sudoku-solver}) boolean-return,
short-circuit-on-success control flow: we need a path through the grid
(mark-before-recurse, unmark-after-return, exactly as in the maze), but
we only need to know \emph{whether one exists} matching the word, so the
search should stop the instant a complete match is found, exactly as
Sudoku stops the instant a full board is completed.
\end{patternbox}

\subsection*{Understanding the problem carefully}

Try starting the search from every cell that matches $\textit{word}$'s
first character (since the path could begin anywhere in the grid).
From a starting cell, the path must match $\textit{word}$ character by
character, moving to an adjacent, not-yet-used cell at each step; the
search succeeds the moment all of $\textit{word}$'s characters have been
matched, and fails along a particular path the moment a required
character cannot be matched by any available neighbor.

\begin{ideabox}[Boolean Grid Backtracking, with In-Place Visited Marking]
$\textsc{Solve}(\textit{row}, \textit{col}, \textit{index})$: if
$\textit{index} = |\textit{word}|$: return \texttt{true} (every
character matched). If out of bounds, or
$\textit{grid}[\textit{row}][\textit{col}] \ne \textit{word}[\textit{index}]$,
or this cell is already used in the current path: return
\texttt{false}. Otherwise: temporarily mark this cell as used (a common
space-saving trick: overwrite the grid cell itself with a sentinel
character unlikely to match anything, rather than maintaining a separate
visited grid). Try all four neighboring directions, recursing with
$\textit{index}+1$; if \emph{any} direction returns \texttt{true}: restore
the cell and return \texttt{true} immediately. If none do: restore the
cell and return \texttt{false}.
\end{ideabox}

\subsection*{Why overwriting the cell in place is a safe substitute for a separate visited array}

Since the grid cell's original character is only needed for the
\emph{duration} of exploring paths through it, and we always restore it
before returning control to the caller (whether the exploration
succeeded or failed), no information is permanently lost -- by the time
any sibling branch or ancestor call inspects this cell again, it has
already been restored to its original character. This avoids allocating
an entire second $m \times n$ boolean grid purely to track visited
cells, at the cost of briefly mutating the input grid during the search
(acceptable as long as the grid is restored to its original state by
the time the overall function returns, which the undo step guarantees).

\subsection*{Algorithm}

\begin{enumerate}
  \item For every cell $(\textit{row}, \textit{col})$ with
        $\textit{grid}[\textit{row}][\textit{col}] = \textit{word}[0]$:
        if $\textsc{Solve}(\textit{row}, \textit{col}, 0)$ returns
        \texttt{true}: the word exists; return \texttt{true}.
  \item $\textsc{Solve}(\textit{row}, \textit{col}, \textit{index})$:
  \begin{enumerate}
    \item If $\textit{index} = |\textit{word}|$: return \texttt{true}.
    \item If out of bounds, or character mismatch, or this exact cell
          is the sentinel (already in use): return \texttt{false}.
    \item Save the original character; overwrite the cell with a
          sentinel.
    \item For each of the $4$ directions: if
          $\textsc{Solve}(\textit{newRow}, \textit{newCol},
          \textit{index}+1)$ returns \texttt{true}: restore the cell;
          return \texttt{true}.
    \item Restore the cell; return \texttt{false}.
  \end{enumerate}
  \item If no starting cell succeeds: return \texttt{false}.
\end{enumerate}

\subsection*{Dry run}

\dryrun{Take the grid
$\begin{bmatrix} A&B&C&E \\ S&F&C&S \\ A&D&E&E \end{bmatrix}$,
$\textit{word} = \texttt{"SEE"}$.}

\begin{itemize}
  \item Starting cells matching \texttt{'S'}: $(1,0)$ and $(1,3)$.
  \item Try $(1,0)$: index $0$ matches \texttt{'S'}. Mark used. Try
        neighbors for index $1$ (\texttt{'E'}): up $(0,0)=\texttt{A}$,
        no; down $(2,0)=\texttt{A}$, no; left out of bounds; right
        $(1,1)=\texttt{F}$, no. All fail. Restore $(1,0)$. Return
        \texttt{false}.
  \item Try $(1,3)$: index $0$ matches \texttt{'S'}. Mark used. Try
        neighbors for index $1$ (\texttt{'E'}): up $(0,3)=\texttt{E}$,
        match! Mark used. Try neighbors for index $2$ (\texttt{'E'}):
        up out of bounds; down (back to $(1,3)$, sentinel, fail); left
        $(0,2)=\texttt{C}$, no; right out of bounds. All fail from
        $(0,3)$. Restore $(0,3)$. Try down from $(1,3)$ for index $1$:
        $(2,3)=\texttt{E}$, match! Mark used. Try neighbors for index
        $2$: up (back to $(1,3)$, sentinel, fail); left
        $(2,2)=\texttt{E}$, match! index $2=3=|\textit{word}|$? Index
        becomes $3$ which equals $|\textit{word}|=3$: return
        \texttt{true} immediately.
  \end{itemize}

The word \texttt{"SEE"} is found via the path $(1,3) \to (2,3) \to
(2,2)$.

\subsection*{C++ implementation}

\begin{lstlisting}
#include <bits/stdc++.h>
using namespace std;

bool solve(int row, int col, int index, vector<vector<char>>& grid,
           string& word) {
    if (index == (int)word.size()) return true;

    int m = grid.size(), n = grid[0].size();
    if (row < 0 || row >= m || col < 0 || col >= n) return false;
    if (grid[row][col] != word[index]) return false;

    char original = grid[row][col];
    grid[row][col] = '#'; // sentinel: marks this cell as in use

    int dr[] = {1, -1, 0, 0};
    int dc[] = {0, 0, 1, -1};
    bool found = false;
    for (int d = 0; d < 4 && !found; d++) {
        found = solve(row + dr[d], col + dc[d], index + 1, grid, word);
    }

    grid[row][col] = original; // restore
    return found;
}

bool wordSearch(vector<vector<char>>& grid, string word) {
    int m = grid.size(), n = grid[0].size();
    for (int r = 0; r < m; r++) {
        for (int c = 0; c < n; c++) {
            if (grid[r][c] == word[0] && solve(r, c, 0, grid, word)) {
                return true;
            }
        }
    }
    return false;
}

int main() {
    vector<vector<char>> grid = {
        {'A','B','C','E'},
        {'S','F','C','S'},
        {'A','D','E','E'}
    };
    cout << (wordSearch(grid, "SEE") ? "true" : "false") << endl;
    return 0;
}
\end{lstlisting}

\begin{complexitybox}
\textbf{Time Complexity.} $O(mn \cdot 4^L)$, where $L$ is the word's
length -- $mn$ possible starting cells, each potentially exploring up to
$4$ directions at every one of $L$ steps.
\textbf{Space Complexity.} $O(L)$ for the recursion depth (no separate
visited grid needed, thanks to the in-place sentinel trick).
\end{complexitybox}

\begin{takeawaybox}
When a grid-backtracking problem only needs a single existence answer,
combine Rat in a Maze's mark/unmark mechanics with Sudoku's
short-circuit-on-success boolean propagation -- and consider whether the
``visited'' marking can be folded directly into the grid itself (via a
temporary sentinel overwrite) rather than allocating a separate
tracking structure, whenever the grid's cells are not otherwise needed
during the brief window before they are restored.
\end{takeawaybox}

\section{Palindrome Partitioning}
\label{sec:rec-palindrome-partitioning}

\problemstatement{Given a string $s$, partition it into one or more
contiguous substrings such that \textbf{every} substring in the
partition is itself a palindrome. Return every such partitioning.}

\begin{patternbox}
\emph{``Partition a sequence into pieces, each satisfying some
property''} introduces a new choice shape: instead of choosing
\emph{which elements} to include (as in subsets/combinations) or
\emph{which direction} to move (as in grid backtracking), the choice here
is \textbf{how far the next piece should extend} -- a loop over every
possible end position for the current piece, starting from wherever the
previous piece left off.
\end{patternbox}

\subsection*{Understanding the problem carefully}

Starting from some position $\textit{start}$ in the string, the next
piece of the partition can end anywhere from $\textit{start}$ itself up
to the end of the string -- but only positions where the resulting
substring $s[\textit{start}..\textit{end}]$ is a palindrome are valid
choices. For each valid choice of where this piece ends, the rest of the
string (from $\textit{end}+1$ onward) must be recursively partitioned
the same way; reaching the very end of the string with no characters
left to partition marks one complete, valid partitioning.

\begin{ideabox}[Loop Over Every Valid Palindromic Prefix from the Current Start]
$\textsc{Solve}(\textit{start}, \textit{current})$: if
$\textit{start} = |s|$: record a copy of $\textit{current}$; return.
Otherwise: for $\textit{end}$ from $\textit{start}$ to $|s|-1$: if
$s[\textit{start}..\textit{end}]$ is a palindrome: append this
substring to $\textit{current}$; recurse
$\textsc{Solve}(\textit{end}+1, \textit{current})$; remove the
substring (undo).
\end{ideabox}

\subsection*{Why looping over end positions correctly enumerates every partitioning}

Every valid partitioning of $s[\textit{start}..]$ is uniquely determined
by where its \emph{first} piece ends -- once that is fixed (to some
valid palindromic substring), the rest of the partitioning is exactly a
partitioning of whatever remains, which is the same sub-problem at a
smaller starting position. By trying \emph{every} valid place the first
piece could end, and recursively solving the remainder for each choice,
we generate every possible partitioning exactly once, with no overlap
between different choices of where the first piece ends (since two
different end positions for the first piece necessarily produce
different first pieces, hence different overall partitionings).

\subsection*{Algorithm}

\begin{enumerate}
  \item $\textsc{Solve}(\textit{start}, \textit{current})$:
  \begin{enumerate}
    \item If $\textit{start} = |s|$: record $\textit{current}$; return.
    \item For $\textit{end}$ from $\textit{start}$ to $|s|-1$:
    \begin{itemize}
      \item If $s[\textit{start}..\textit{end}]$ is a palindrome:
            append it to $\textit{current}$; recurse
            $\textsc{Solve}(\textit{end}+1, \textit{current})$; remove
            it (undo).
    \end{itemize}
  \end{enumerate}
  \item Call $\textsc{Solve}(0, [\,])$.
\end{enumerate}

\subsection*{Dry run}

\dryrun{Take $s = \texttt{"aab"}$.}

\begin{itemize}
  \item $\textsc{Solve}(0,[])$: $\textit{end}=0$:
        $s[0..0]=\texttt{"a"}$, palindrome. Append. Recurse
        $\textsc{Solve}(1,[\texttt{"a"}])$.
  \begin{itemize}
    \item $\textit{end}=1$: $s[1..1]=\texttt{"a"}$, palindrome.
          Append. Recurse $\textsc{Solve}(2,[\texttt{"a"},\texttt{"a"}])$.
    \begin{itemize}
      \item $\textit{end}=2$: $s[2..2]=\texttt{"b"}$, palindrome.
            Append. Recurse $\textsc{Solve}(3,[\texttt{"a"},\texttt{"a"},\texttt{"b"}])$:
            $\textit{start}=3=|s|$, record
            $[\texttt{"a"},\texttt{"a"},\texttt{"b"}]$.
    \end{itemize}
    \item Undo. $\textit{end}=2$ at the $[\texttt{"a"}]$ level:
          $s[1..2]=\texttt{"ab"}$, not a palindrome. Skip.
  \end{itemize}
  \item Undo back to $[]$. $\textit{end}=1$ at the top level:
        $s[0..1]=\texttt{"aa"}$, palindrome. Append. Recurse
        $\textsc{Solve}(2,[\texttt{"aa"}])$.
  \begin{itemize}
    \item $\textit{end}=2$: $s[2..2]=\texttt{"b"}$, palindrome.
          Append. Recurse $\textsc{Solve}(3,[\texttt{"aa"},\texttt{"b"}])$:
          record $[\texttt{"aa"},\texttt{"b"}]$.
  \end{itemize}
  \item $\textit{end}=2$ at the top level: $s[0..2]=\texttt{"aab"}$,
        not a palindrome. Skip.
\end{itemize}

Recorded partitionings: $[\texttt{"a"},\texttt{"a"},\texttt{"b"}]$ and
$[\texttt{"aa"},\texttt{"b"}]$ -- the only two ways to partition
\texttt{"aab"} into all-palindrome pieces.

\subsection*{C++ implementation}

\begin{lstlisting}
#include <bits/stdc++.h>
using namespace std;

bool isPalindrome(string& s, int lo, int hi) {
    while (lo < hi) {
        if (s[lo] != s[hi]) return false;
        lo++; hi--;
    }
    return true;
}

void solve(int start, string& s, vector<string>& current,
           vector<vector<string>>& result) {
    if (start == (int)s.size()) {
        result.push_back(current);
        return;
    }
    for (int end = start; end < (int)s.size(); end++) {
        if (isPalindrome(s, start, end)) {
            current.push_back(s.substr(start, end - start + 1));
            solve(end + 1, s, current, result);
            current.pop_back();
        }
    }
}

vector<vector<string>> palindromePartitioning(string s) {
    vector<vector<string>> result;
    vector<string> current;
    solve(0, s, current, result);
    return result;
}

int main() {
    for (auto& partition : palindromePartitioning("aab")) {
        cout << "[ ";
        for (string& piece : partition) cout << piece << " ";
        cout << "] ";
    }
    cout << endl;
    return 0;
}
\end{lstlisting}

\begin{complexitybox}
\textbf{Time Complexity.} Up to $O(2^n)$ partitionings can exist in the
worst case (e.g.\ a string of all-identical characters, where every
position could start a new piece), and checking each candidate substring
for the palindrome property costs up to $O(n)$, giving an overall worst
case of
\[
  O(n \cdot 2^n).
\]
\textbf{Space Complexity.} $O(n)$ for the recursion depth.
\end{complexitybox}

\begin{takeawaybox}
``Partition a sequence into pieces satisfying a property'' is a third
fundamental backtracking choice shape, alongside ``include or exclude
each element'' and ``move in one of several directions'': loop over
every possible \emph{extent} of the next piece from the current
position, recurse on the remainder for each valid extent. This same
``loop over partition points'' shape reappears, generalized with memoization,
in the Matrix Chain Multiplication / Partition DP family covered later in
the dynamic programming chapters.
\end{takeawaybox}

\section{Word Break}
\label{sec:rec-word-break}

\problemstatement{Given a string $s$ and a dictionary of words
$\textit{wordDict}$, determine whether $s$ can be segmented into a
sequence of one or more dictionary words placed consecutively (words may
be reused any number of times).}

\begin{patternbox}
Structurally, this is Palindrome Partitioning's
(Section~\ref{sec:rec-palindrome-partitioning}) ``loop over the next
piece's extent'' choice shape, combined with Sudoku's and Word Search's
boolean short-circuit-on-success control flow -- but with one crucial
difference that makes this section the deliberate bridge into the
dynamic programming chapters that follow: the sub-problem
``can $s[\textit{start}..]$ be segmented'' depends only on
$\textit{start}$, and the \textbf{same} value of $\textit{start}$ can be
reached through many different earlier choices of piece boundaries --
these are \textbf{overlapping subproblems}, the defining signature of a
problem where memoization will pay off enormously.
\end{patternbox}

\subsection*{Understanding the problem carefully}

From position $\textit{start}$, try every possible end position for the
next word: if $s[\textit{start}..\textit{end}]$ is in the dictionary,
recursively check whether the remainder, $s[\textit{end}+1..]$, can also
be fully segmented. If reaching the very end of the string is possible
through any sequence of such choices, $s$ is breakable. As in Sudoku
(Section~\ref{sec:rec-sudoku-solver}), the moment any choice leads to a
fully successful segmentation of the remainder, we can stop trying
further alternatives and propagate \texttt{true} immediately.

\begin{ideabox}[Loop Over Dictionary-Matching Prefixes, Short-Circuit on Success]
$\textsc{Solve}(\textit{start})$: if $\textit{start} = |s|$: return
\texttt{true} (the whole string has been consumed). Otherwise, for
$\textit{end}$ from $\textit{start}$ to $|s|-1$: if
$s[\textit{start}..\textit{end}]$ is in $\textit{wordDict}$ and
$\textsc{Solve}(\textit{end}+1)$ returns \texttt{true}: return
\texttt{true} immediately. If no choice of $\textit{end}$ leads to
success: return \texttt{false}.
\end{ideabox}

\subsection*{Why this recursion, run naively, repeats the same work exponentially often}

Consider a string with many ways to reach the same intermediate position
$\textit{start}$ -- for instance, $s=\texttt{"aaaaaaaaaaaaaaaaaaaaaaaab"}$
with dictionary $\{\texttt{"a"},\texttt{"aa"},\texttt{"aaa"}\}$: the
position right after consuming, say, the first $10$ characters can be
reached by choosing word lengths $1+1+\cdots$, or $2+2+\cdots$, or
$1+2+1+\cdots$, and many other combinations -- yet
$\textsc{Solve}(10)$ always asks exactly the same question (``can
$s[10..]$ be segmented'') regardless of \emph{which} combination of
earlier word choices got us to position $10$. A naive recursive
implementation recomputes the answer to $\textsc{Solve}(10)$ separately,
from scratch, every single time any path happens to reach it -- and the
number of distinct paths reaching a given position can grow
exponentially with the string's length, even though there are only
$|s|+1$ \emph{distinct} values $\textit{start}$ could ever take.

\subsection*{The bridge to dynamic programming}

This is precisely the situation dynamic programming exists to fix:
\textbf{memoization} caches the result of $\textsc{Solve}(\textit{start})$
the first time it is computed, in an array (or hash map) indexed by
$\textit{start}$, and every subsequent call with the same
$\textit{start}$ simply looks up the cached answer in $O(1)$ instead of
re-deriving it -- collapsing the exponential blowup down to $O(|s|)$
distinct sub-problems, each solved once. The recursion's
\emph{structure} (the recurrence relating $\textsc{Solve}(\textit{start})$
to $\textsc{Solve}(\textit{end}+1)$ for various $\textit{end}$) does not
change at all; only the bookkeeping around it does. This exact
transformation -- recognize overlapping subproblems in an already-correct
recursion, then cache results by the parameters that define each distinct
subproblem -- is the opening idea of the next chapter, where it is
applied systematically across dozens of problems, progressing from this
kind of memoized top-down recursion to fully iterative tabulation and,
where possible, space-optimized tabulation.

\subsection*{Algorithm}

\begin{enumerate}
  \item $\textsc{Solve}(\textit{start})$:
  \begin{enumerate}
    \item If $\textit{start}=|s|$: return \texttt{true}.
    \item For $\textit{end}$ from $\textit{start}$ to $|s|-1$: if
          $s[\textit{start}..\textit{end}] \in \textit{wordDict}$ and
          $\textsc{Solve}(\textit{end}+1)$: return \texttt{true}.
    \item Return \texttt{false}.
  \end{enumerate}
  \item Call $\textsc{Solve}(0)$.
\end{enumerate}

\subsection*{Dry run}

\dryrun{Take $s = \texttt{"leetcode"}$, $\textit{wordDict} =
\{\texttt{"leet"}, \texttt{"code"}\}$.}

\begin{itemize}
  \item $\textsc{Solve}(0)$: try $\textit{end}=0..3$:
        $s[0..3]=\texttt{"leet"}$ is in the dictionary. Check
        $\textsc{Solve}(4)$.
  \item $\textsc{Solve}(4)$: try $\textit{end}=4..7$:
        $s[4..7]=\texttt{"code"}$ is in the dictionary. Check
        $\textsc{Solve}(8)$.
  \item $\textsc{Solve}(8)$: $\textit{start}=8=|s|$. Return
        \texttt{true}.
  \item This propagates straight back: $\textsc{Solve}(4)$ returns
        \texttt{true}; $\textsc{Solve}(0)$ returns \texttt{true}.
\end{itemize}

The string segments as \texttt{"leet"} + \texttt{"code"}.

\subsection*{C++ implementation}

\begin{lstlisting}
#include <bits/stdc++.h>
using namespace std;

bool solve(int start, string& s, unordered_set<string>& wordSet) {
    if (start == (int)s.size()) return true;

    for (int end = start; end < (int)s.size(); end++) {
        string piece = s.substr(start, end - start + 1);
        if (wordSet.count(piece) && solve(end + 1, s, wordSet)) {
            return true;
        }
    }
    return false;
}

bool wordBreak(string s, vector<string>& wordDict) {
    unordered_set<string> wordSet(wordDict.begin(), wordDict.end());
    return solve(0, s, wordSet);
}

int main() {
    vector<string> wordDict = {"leet", "code"};
    cout << (wordBreak("leetcode", wordDict) ? "true" : "false") << endl;
    return 0;
}
\end{lstlisting}

\begin{complexitybox}
\textbf{Time Complexity.} $O(2^n)$ in the worst case for this naive
(unmemoized) version, due to the overlapping-subproblem blowup described
above. With memoization keyed by $\textit{start}$ (an array or hash map
of size $O(n)$, each computed once in $O(n)$ work to try all
$\textit{end}$ positions), this drops to $O(n^2)$.
\textbf{Space Complexity.} $O(n)$ for the recursion depth, plus
$O(n)$ for the memoization cache once added.
\end{complexitybox}

\begin{takeawaybox}
This closing problem of the Recursion chapter is a deliberate hinge:
everything about its recursive structure is already correct and familiar
from this chapter's earlier sections, but its overlapping subproblems
expose recursion's biggest weakness -- redundant recomputation -- setting
up exactly the problem that memoization and tabulation, covered next, are
designed to solve. Carry this recursion forward unchanged; only the
bookkeeping around it is about to be upgraded.
\end{takeawaybox}

